\studySubsection{calc-01-function-limit-continuity}{第 1 讲 函数极限与连续}

\begin{knowledgeBox}
教材页码：第 1 页起。

本讲用于整理函数、函数极限、无穷小比较、连续性与间断点等内容。新增题目时重点沉淀极限入口、等价替换条件、单侧极限和连续性判定。
\end{knowledgeBox}

\knowledgeAnchor[MATH1-KN-CALC-0059]{七种未定式的统一化归}
\noindent{\footnotesize\textbf{索引：} \knowledgeIndexRef[MATH1-KN-CALC-0059]{七种未定式的统一化归}}\par

\begin{knowledgeBox}
\textbf{来源定位：}

张宇基础 30 讲（高数），PDF 第 60 页至第 70 页上半页（教材页码第 54 页至第 64 页上半页）“七种未定式的计算”，由用户指定。

\textbf{什么叫未定式：}

函数极限中常见的七种未定式是
\[
\frac{0}{0},\qquad
\frac{\infty}{\infty},\qquad
0\cdot\infty,\qquad
\infty-\infty,\qquad
\infty^0,\qquad
0^0,\qquad
1^\infty.
\]
这里的 \(0\) 和 \(\infty\) 表示各部分的极限状态，不是把两个数直接代入运算。“未定”只表示这些状态本身不能唯一决定原极限；原极限可能是有限数、\(0\)、无穷，也可能不存在。

\textbf{统一路线：}

七种形式并不需要背七套互不相干的算法。真正直接处理的只有两种商式
\[
\frac{0}{0},\qquad \frac{\infty}{\infty}.
\]
其余五种先化归：积式改写成商式，差式通分或有理化，幂式先取对数。完整顺序是
\[
\text{代入识别}
\longrightarrow
\text{先做代数化简}
\longrightarrow
\text{化为商式}
\longrightarrow
\text{计算极限}.
\]

\textbf{一、\(\frac{0}{0}\) 型：先消掉共同趋零的原因。}

常用入口依次是因式分解与约分、有理化、等价无穷小、泰勒展开、洛必达法则。例如
\[
\lim_{x\to0}\frac{e^x-1-x}{x^2}
=\lim_{x\to0}\frac{e^x-1}{2x}
=\lim_{x\to0}\frac{e^x}{2}
=\frac12,
\]
其中连续两次使用洛必达法则，每次使用前都仍是 \(\frac00\) 型。

\textbf{二、\(\frac{\infty}{\infty}\) 型：比较分子、分母的增长阶。}

多项式或根式优先同除最高阶量；也可以使用等价量、泰勒展开或洛必达法则。例如
\[
\lim_{x\to+\infty}\frac{3x^2+x}{2x^2-1}
=\lim_{x\to+\infty}\frac{3+\frac1x}{2-\frac1{x^2}}
=\frac32.
\]

\textbf{三、\(0\cdot\infty\) 型：把更容易取倒数的一项放到分母。}

若 \(f(x)\to0\)、\(g(x)\to\infty\)，则可以选择
\[
f(x)g(x)=\frac{f(x)}{1/g(x)}
\quad\text{或}\quad
f(x)g(x)=\frac{g(x)}{1/f(x)},
\]
分别化为 \(\frac00\) 型或 \(\frac{\infty}{\infty}\) 型。选择导数更简单的一种。例如
\[
\lim_{x\to0^+}x\ln x
=\lim_{x\to0^+}\frac{\ln x}{1/x}
=\lim_{x\to0^+}\frac{1/x}{-1/x^2}
=\lim_{x\to0^+}(-x)=0.
\]

\textbf{四、\(\infty-\infty\) 型：把相消关系显出来。}

分式差优先通分，根式差优先乘共轭式，对数差优先合并成一个对数。例如
\[
\begin{aligned}
\lim_{x\to+\infty}\left(\sqrt{x^2+x}-x\right)
&=\lim_{x\to+\infty}
\frac{x}{\sqrt{x^2+x}+x}\\
&=\lim_{x\to+\infty}
\frac{1}{\sqrt{1+1/x}+1}
=\frac12.
\end{aligned}
\]
洛必达法则不能直接作用于“差”，必须先把它化成一个商。

\textbf{五、六、七：三种幂指未定式统一取对数。}

设在所讨论的去心邻域内 \(f(x)>0\)，并令
\[
y=f(x)^{g(x)}.
\]
则
\[
\ln y=g(x)\ln f(x).
\]
无论原式是 \(1^\infty\)、\(0^0\) 还是 \(\infty^0\)，右端都会化为 \(0\cdot\infty\) 型；先求
\[
L=\lim g(x)\ln f(x),
\]
再还原
\[
\lim f(x)^{g(x)}=e^L.
\]

三种典型例子分别是
\[
\lim_{x\to+\infty}\left(1+\frac1x\right)^x=e,
\]
因为 \(x\ln(1+1/x)\to1\)；
\[
\lim_{x\to0^+}x^x=1,
\]
因为 \(x\ln x\to0\)；以及
\[
\lim_{x\to+\infty}x^{1/x}=1,
\]
因为 \(\frac{\ln x}{x}\to0\)。

对 \(1^\infty\) 型还有常用速算式：若 \(u(x)\to0\)、\(v(x)\to\infty\)、\(1+u(x)>0\)，且
\[
\lim u(x)v(x)=A,
\]
则由 \(\ln(1+u)\sim u\) 得
\[
\lim[1+u(x)]^{v(x)}=e^A.
\]
因此 \(1^\infty\) 型的答案不一定是 \(e\)，而是由指数中的极限 \(A\) 决定。

\textbf{后续页的深化：先化简，再判型，最后选法。}

教材后续例题反复强调的不是“看到极限就求导”，而是下面三层决策：
\[
\text{化简结构}
\longrightarrow
\text{判断真正的未定式}
\longrightarrow
\text{在等价替换、泰勒、洛必达和夹逼中选最短路线}.
\]
化简包括约分、提取极限不为零的因子、通分、有理化、合并、变量代换，以及把分子分母同时除以合适的最高次或最低次幂。很多题若不先化简，后续计算会急剧膨胀。

\textbf{等价无穷小只保证乘除替换，不保证加减替换。}

若 \(u(x)\sim v(x)\)，则它们的比值趋于 \(1\)，所以在乘积或商中通常可以替换；但在差
\[
u(x)-v(x)
\]
中，两项的主部可能恰好抵消，此时必须保留更高阶项。例如只知道
\[
\ln(1+x)\sim x
\]
无法算出 \(x-\ln(1+x)\) 的主部。需要升到二阶：
\[
\ln(1+x)=x-\frac{x^2}{2}+o(x^2),
\qquad
e^x=1+x+\frac{x^2}{2}+o(x^2).
\]
于是得到两组高频结论
\[
x-\ln(1+x)\sim\frac{x^2}{2},
\qquad
e^x-1-x\sim\frac{x^2}{2}.
\]
同理，
\[
\frac{\ln(1+x)}{x}
=1-\frac{x}{2}+o(x),
\]
所以
\[
(1+x)^{1/x}
=e^{\,1-\frac{x}{2}+o(x)}
=e-\frac{e}{2}x+o(x).
\]
这说明当一阶主部发生相消时，题目真正考的是“找到第一个不为零的高阶项”。

\textbf{“抓大头”前必须检查系数不会消失。}

若
\[
P_n(x)=a_nx^n+\cdots,\qquad
Q_m(x)=b_mx^m+\cdots,\qquad a_nb_m\ne0,
\]
则当 \(x\to\infty\) 时，
\[
\frac{P_n(x)}{Q_m(x)}
\sim\frac{a_n}{b_m}x^{n-m}.
\]
因此 \(n=m\) 时取最高次项系数之比，\(n>m\) 时再结合符号判断无穷方向，\(n<m\) 时极限为 \(0\)。当 \(x\to0\) 时逻辑相反，应抓最低次非零项。

含根式且 \(x\to-\infty\) 时还必须处理
\[
\sqrt{x^2}=|x|=-x,
\]
不能把它机械写成 \(x\)。

若“最高次项”的系数依赖参数或自变量，并且可能等于零，则不能统一抓大头，必须先按该系数是否为零分情况。例如
\[
F(x)=\lim_{n\to\infty}
\frac{A(x)+nB(x)}{C(x)+nD(x)}
\]
只有在 \(D(x)\ne0\) 时才可直接抓 \(n\) 的系数得到 \(B(x)/D(x)\)；当 \(B(x)\) 与 \(D(x)\) 同时为零时，含 \(n\) 的项全部消失，必须回到剩余项重新计算。这正是极限定义函数经常产生分段结果的原因。

\textbf{已知一个极限反求未知函数：脱帽法。}

若
\[
\lim_{x\to a}F(x)=L,
\]
则可以写成
\[
F(x)=L+\alpha(x),\qquad \alpha(x)\to0.
\]
这不是近似猜测，而是极限定义的直接改写。若题设给出
\[
\frac{A(x)+B(x)f(x)}{\varphi(x)}\to L,
\]
则在 \(B(x)\ne0\) 的去心邻域内可写成
\[
A(x)+B(x)f(x)
=\bigl[L+\alpha(x)\bigr]\varphi(x),
\qquad \alpha(x)\to0,
\]
从而解出 \(f(x)\)，再代入目标极限。若已知极限为 \(0\)，就得到分子是 \(o(\varphi(x))\)。这种方法特别适合“已知一极限，求另一极限”的题；泰勒展开则负责给出 \(A(x)\) 的精确阶数。

\textbf{不连续或取整结构优先考虑夹逼。}

取整函数本身在相关点附近可能没有极限，但它与连续小量相乘后可以收敛。利用
\[
\lfloor t\rfloor=t+\theta(t),
\qquad -1<\theta(t)\le0,
\]
可得
\[
x\left\lfloor\frac{a}{x}\right\rfloor
=a+x\theta\left(\frac{a}{x}\right)\to a
\qquad (x\to0).
\]
这里真正起作用的是取整误差始终有界，而 \(x\to0\) 会把该误差压到零。遇到取整、有限振荡或其他有界扰动时，先寻找
\[
\text{趋零因子}\times\text{有界误差},
\]
往往比强行使用洛必达更自然。

\textbf{增长阶：任意正幂都能压住对数。}

对任意 \(\alpha>0\)、\(\beta>0\)，有
\[
\lim_{x\to0^+}x^\alpha|\ln x|^\beta=0,
\qquad
\lim_{x\to+\infty}\frac{(\ln x)^\beta}{x^\alpha}=0.
\]
所以 \(x\ln x\to0\) 不是孤立公式，而是“幂函数比对数函数变化更快”的一个特例。它同时服务于 \(0\cdot\infty\)、\(0^0\) 和 \(\infty^0\) 型。

\textbf{\(\infty-\infty\) 型若没有现成分母，就主动制造小变量。}

分式差先通分，根式差先有理化；若原式没有分母，可以提取公因子，再令 \(u=1/x\) 把无穷远问题搬到零点附近。例如
\[
\begin{aligned}
\lim_{x\to+\infty}
\left[x^2\left(e^{1/x}-1\right)-x\right]
&=\lim_{u\to0^+}
\frac{e^u-1-u}{u^2}\\
&=\frac12.
\end{aligned}
\]
这样原来的 \(\infty-\infty\) 就被化为标准的二阶 \(\frac00\) 型。

\textbf{泰勒展开的本质是系数匹配。}

若要求
\[
F(x)-P_k(x)=o(x^k),
\]
则多项式 \(P_k(x)\) 必须与 \(F(x)\) 的泰勒展开从常数项到 \(x^k\) 项逐项一致。例如
\[
e^x-\left(1+bx+ax^2\right)=o(x^2)
\]
迫使
\[
b=1,\qquad a=\frac12.
\]
若函数具有奇偶性，还可以先删除不可能出现的系数：奇函数的泰勒多项式只含奇次幂，偶函数只含偶次幂。例如
\[
\frac{\sin x}{1+x^2}
=\left(x-\frac{x^3}{6}+o(x^3)\right)
\left(1-x^2+o(x^2)\right)
=x-\frac76x^3+o(x^3).
\]
泰勒特别适合主部相消、要求高阶极限、反求参数或求局部多项式；这几类题只用一阶等价无穷小通常不够。

\textbf{工具选择顺序：}

\begin{enumerate}
  \item 能通过恒等变形、约分、通分、有理化或抓主项直接结束，就不用重工具。
  \item 纯乘除结构优先等价无穷小；出现加减相消时，改用足够阶的泰勒展开。
  \item 已经化为 \(\frac00\) 或 \(\frac{\infty}{\infty}\)，且求导后明显更简单时，再用洛必达法则。
  \item 含取整、振荡或有界扰动时，优先把误差隔离出来并使用夹逼。
  \item 含参数或由极限定义函数时，先检查被抓取的主项系数是否会变成零；会变零就分情况。
\end{enumerate}

\textbf{洛必达法则的使用边界：}

洛必达法则只直接用于 \(\frac00\) 型或 \(\frac{\infty}{\infty}\) 型。使用时要保证分子、分母在相应去心邻域内可导，分母导数不为零，并且导数之比的极限存在或为无穷。每使用一次后都要重新判型；若已经不再是这两种商式，不能机械地继续求导。

\textbf{考场检查清单：}

\begin{enumerate}
  \item 先代入判型，不见未定式就不要启动洛必达法则。
  \item 先约分、提非零因子、有理化、通分或同除最高阶量；化简往往比求导更快。
  \item \(0\cdot\infty\) 改商，\(\infty-\infty\) 合并成商，三种幂式取对数。
  \item 幂指函数取对数前检查底数最终为正；底数可能为负时，原式对一般实指数甚至未必有定义。
  \item 得到答案后检查单侧、符号和定义域，尤其注意 \(x\to0^+\) 与 \(x\to0^-\) 可能不同。
\end{enumerate}

\textbf{主动回忆：}

合上答案后，口述三句话：“商式直接处理；积和差先化商；幂式先取对数。”再回答五个问题：为什么 \(x-\ln(1+x)\) 不能只用一阶等价替换；抓大头前要检查什么；脱帽法把已知极限改写成什么；取整函数题为什么常用夹逼；高阶无穷小条件怎样转化为泰勒系数匹配。
\end{knowledgeBox}

\studySubsection{calc-01-inverse-function}{反函数的定义、存在条件与图像关系}
\knowledgeAnchor[MATH1-KN-CALC-0008]{反函数}
\legacyKnowledgeAnchor{inverse-function}
\noindent{\footnotesize\textbf{索引：} \knowledgeIndexRef[MATH1-KN-CALC-0008]{反函数}}\par

\begin{knowledgeBox}
\textbf{概念定位：}

反函数属于函数基础中的“对应关系”问题。本质不是计算技巧，而是判断一个函数值能否唯一反推出原来的自变量。后续求反函数、反函数求导、复合函数化简、图像对称和定义域值域问题都会用到它。

\textbf{定义：}

设函数 \(y=f(x)\) 的定义域为 \(D\)，值域为 \(R\)。如果对每一个 \(y\in R\)，在 \(D\) 中都存在唯一的 \(x\)，使得
\[
y=f(x),
\]
则可以由 \(y\) 唯一确定 \(x\)。这时定义新函数
\[
x=\varphi(y),
\]
称它为 \(y=f(x)\) 的反函数，通常记作
\[
x=f^{-1}(y).
\]
反函数的定义域是原函数的值域 \(R\)，反函数的值域是原函数的定义域 \(D\)。

\textbf{直觉解释：}

原函数是“输入 \(x\)，得到 \(y\)”；反函数是“已知输出 \(y\)，倒推输入 \(x\)”。所以反函数存在的核心要求是：同一个输出不能来自两个不同输入。若存在 \(x_1\ne x_2\)，但 \(f(x_1)=f(x_2)\)，那么只知道 \(y\) 时就无法判断原来的输入到底是 \(x_1\) 还是 \(x_2\)，反函数就不存在。

\textbf{存在条件：}

函数有反函数的充要条件是它在定义域上是一一对应，即
\[
x_1\ne x_2\quad\Longrightarrow\quad f(x_1)\ne f(x_2).
\]
严格单调函数一定一一对应，所以严格单调是反函数存在的充分条件。但它不是必要条件：有反函数的函数不一定在整个定义域上单调。

\textbf{典型例子一：限制定义域后才有反函数}

函数 \(y=x^2\) 若定义在全体实数上，则 \(f(1)=f(-1)=1\)，同一个函数值 \(1\) 对应两个不同自变量，因此没有反函数。

如果把定义域限制为 \([0,+\infty)\)，则 \(y=x^2\) 严格递增，且值域为 \([0,+\infty)\)，于是有反函数
\[
x=\sqrt{y}.
\]
按通常习惯把反函数也写成以 \(x\) 为自变量，就得到
\[
y=\sqrt{x},\qquad x\ge 0.
\]
这说明：很多函数不是天然没有反函数，而是需要先选定一个能保证一一对应的定义域分支。

\textbf{典型例子二：有反函数不一定单调}

设
\[
f(x)=
\begin{cases}
x, & x\ge 0,\\
\dfrac{1}{x}, & x<0.
\end{cases}
\]
当 \(x\ge 0\) 时，函数值落在 \([0,+\infty)\)；当 \(x<0\) 时，函数值落在 \((-\infty,0)\)。两段的值域不重叠，且每段内部都不会出现同一个函数值对应两个自变量，所以它是一一对应的，有反函数。事实上它的反函数就是它本身。

但这个函数在整个定义域上不是单调函数。例如在负半轴上，\(1/x\) 是递减的，而在非负半轴上，\(x\) 是递增的；跨过 \(0\) 时又不是一个统一的单调趋势。因此“严格单调必有反函数”可以用，“有反函数必严格单调”不能用。

\textbf{图像关系：}

若把
\[
x=f^{-1}(y)
\]
和
\[
y=f(x)
\]
画在同一坐标系中，它们其实表示同一批点，只是从不同方向描述同一个对应关系。

若把反函数改写成标准函数形式 \(y=f^{-1}(x)\)，相当于把原关系中的 \(x\) 与 \(y\) 互换。因此 \(y=f(x)\) 与 \(y=f^{-1}(x)\) 的图像关于直线
\[
y=x
\]
对称。

\textbf{判别方法：}

\begin{itemize}
  \item 铅直线判定函数：任意一条铅直线与图像至多一个交点，图像才表示 \(y\) 是 \(x\) 的函数。
  \item 水平线判定反函数：在已经是函数的前提下，任意一条水平线与图像至多一个交点，函数才有反函数。
\end{itemize}

\textbf{求反函数五步法：}

\begin{enumerate}
  \item 写出原函数 \(y=f(x)\)，并明确原定义域 \(D\)。
  \item 由方程 \(y=f(x)\) 解出 \(x\)，得到 \(x=\varphi(y)\)。
  \item 检查解出的 \(x\) 是否唯一；若不唯一，要限制定义域或分支。
  \item 交换字母 \(x,y\)，写成 \(y=f^{-1}(x)\)。
  \item 写明反函数的定义域和值域：反函数定义域为原函数值域，反函数值域为原函数定义域。
\end{enumerate}

\textbf{重要关系：}

若 \(f:D\to R\) 一一对应，且反函数为 \(f^{-1}:R\to D\)，则
\[
f^{-1}(f(x))=x,\qquad x\in D,
\]
\[
f(f^{-1}(y))=y,\qquad y\in R.
\]
写这些恒等式时必须带着变量范围理解：第一个式子里的 \(x\) 来自原函数定义域，第二个式子里的 \(y\) 来自原函数值域。

\textbf{常见误区：}

\begin{itemize}
  \item \(f^{-1}(x)\) 表示反函数，不表示 \(\dfrac{1}{f(x)}\)。
  \item 只说“单调所以有反函数”时，要确认是在整个给定定义域上严格单调。
  \item 求反函数后忘记写定义域和值域，会在根式、对数和分段函数中丢分。
  \item 图像对称关系只适用于 \(y=f(x)\) 与 \(y=f^{-1}(x)\)，不是 \(x=f^{-1}(y)\) 与 \(y=f(x)\) 两条不同曲线关于 \(y=x\) 对称。
\end{itemize}
\end{knowledgeBox}

\begin{mistakeBox}
\textbf{满分复盘提醒：}

\begin{itemize}
  \item 判断反函数存在时，第一句话应落到“一一对应”，不要只背“单调”。
  \item 看到 \(x^2\)、三角函数、根式函数时，优先检查是否需要限制定义域分支。
  \item 看到 \(f^{-1}(f(x))\) 或 \(f(f^{-1}(x))\) 时，先检查内层函数值是否落在外层函数定义域内。
  \item 画图题中要区分铅直线判定“是不是函数”和水平线判定“有没有反函数”。
\end{itemize}
\end{mistakeBox}

\studySubsection{calc-01-inverse-hyperbolic-sine}{反双曲正弦型函数的反函数}

\problemAnchor{MATH1-CALC-0003}
\begin{problemBox}
\begin{problemMeta}
\textbf{题目编号：} \problemRef{MATH1-CALC-0003} \quad
\textbf{索引：} \problemIndexRef{MATH1-CALC-0003} \quad
\textbf{来源：} 用户图片 / 教材例 1.3

\textbf{题型：} 求反函数表达式与定义域 \quad
\textbf{难度：} 基础

\textbf{知识点：} \knowledgeRef[MATH1-KN-CALC-0008]{反函数}，对数性质，有理化，双曲函数
\end{problemMeta}

\textbf{题目：}

求函数
\[
y=f(x)=\ln\left(x+\sqrt{x^2+1}\right)
\]
的反函数 \(f^{-1}(x)\) 的表达式及其定义域。
\end{problemBox}

\begin{solutionBox}
\textbf{思路入口：}

本题表面上是“从
\[
y=\ln\left(x+\sqrt{x^2+1}\right)
\]
中解出 \(x\)”，但若直接去解含根式方程，容易把式子平方得很繁。更好的入口是注意到
\[
\left(x+\sqrt{x^2+1}\right)\left(\sqrt{x^2+1}-x\right)=1.
\]
也就是说，\(x+\sqrt{x^2+1}\) 的倒数正好是 \(\sqrt{x^2+1}-x\)。这会把对数式变成一对互补关系：
\[
e^y=x+\sqrt{x^2+1},\qquad e^{-y}=\sqrt{x^2+1}-x.
\]
从思维上看，这仍然是方程组思维：把 \(x\) 和 \(\sqrt{x^2+1}\) 暂时看成两个需要消元的对象，通过两条同源方程联立，最后用相减消去根式。

\textbf{详细解答：}

由题设
\[
y=\ln\left(x+\sqrt{x^2+1}\right).
\]
因为 \(\sqrt{x^2+1}>|x|\)，所以
\[
x+\sqrt{x^2+1}>0,
\]
对数对任意实数 \(x\) 都有意义。两边取指数，得
\[
e^y=x+\sqrt{x^2+1}. \tag{1}
\]
再看 \(-y\)。由对数性质
\[
-y=-\ln\left(x+\sqrt{x^2+1}\right)
=\ln\frac{1}{x+\sqrt{x^2+1}}.
\]
对分母有理化：
\[
\frac{1}{x+\sqrt{x^2+1}}
=\frac{\sqrt{x^2+1}-x}{\left(\sqrt{x^2+1}+x\right)\left(\sqrt{x^2+1}-x\right)}
=\sqrt{x^2+1}-x.
\]
所以
\[
-y=\ln\left(\sqrt{x^2+1}-x\right),
\]
进而
\[
e^{-y}=\sqrt{x^2+1}-x. \tag{2}
\]
此时 \((1),(2)\) 可以看成关于 \(x\) 与 \(\sqrt{x^2+1}\) 的二元一次方程组。用 \((1)-(2)\)，根式项抵消：
\[
e^y-e^{-y}=2x.
\]
因此
\[
x=\frac{1}{2}\left(e^y-e^{-y}\right).
\]
这已经把原来的自变量 \(x\) 表示成了函数值 \(y\) 的函数，即
\[
x=f^{-1}(y)=\frac{1}{2}\left(e^y-e^{-y}\right).
\]
按通常书写习惯，将自变量记回 \(x\)，得到
\[
f^{-1}(x)=\frac{1}{2}\left(e^x-e^{-x}\right).
\]

\textbf{定义域与值域：}

原函数的定义域为
\[
(-\infty,+\infty).
\]
由
\[
\eta=\frac{1}{2}\left(e^y-e^{-y}\right)
\]
可知，对任意 \(y\in\mathbb{R}\)，都能找到唯一实数 \(\eta\) 使 \(f(\eta)=y\)。所以原函数值域为 \(\mathbb{R}\)，反函数的定义域为
\[
(-\infty,+\infty).
\]
对应地，反函数的值域等于原函数定义域，也是
\[
(-\infty,+\infty).
\]
于是最终答案为
\[
\boxed{f^{-1}(x)=\frac{1}{2}\left(e^x-e^{-x}\right),\qquad x\in(-\infty,+\infty)}.
\]

\textbf{为什么它一定有反函数：}

本题不只要会算，还要知道为什么可以谈反函数。函数
\[
f(x)=\ln\left(x+\sqrt{x^2+1}\right)
\]
是严格递增函数。可以从反函数表达式看出，\(x=\frac12(e^y-e^{-y})\) 对每个 \(y\in\mathbb{R}\) 都给出唯一的 \(x\)；也可以在后续学到导数后验证
\[
f'(x)=\frac{1}{\sqrt{x^2+1}}>0.
\]
因此原函数一一对应，有反函数。

\textbf{考场满分写法：}

必须写：
\[
e^y=x+\sqrt{x^2+1},\qquad
e^{-y}=\sqrt{x^2+1}-x.
\]
其中第二式要通过倒数有理化说明：
\[
\frac{1}{x+\sqrt{x^2+1}}=\sqrt{x^2+1}-x.
\]
然后相减得
\[
x=\frac12(e^y-e^{-y}).
\]
最后交换 \(x,y\)，并写出
\[
f^{-1}(x)=\frac12(e^x-e^{-x}),\qquad x\in\mathbb{R}.
\]
可以略写的是原函数单调性的详细证明，但不能漏掉反函数定义域。

\textbf{答案：}

\[
\boxed{f^{-1}(x)=\frac{1}{2}\left(e^x-e^{-x}\right),\quad x\in\mathbb{R}}.
\]
\end{solutionBox}

\begin{knowledgeBox}
\textbf{核心知识点：}

函数
\[
\ln\left(x+\sqrt{x^2+1}\right)
\]
常记为反双曲正弦函数 \(\operatorname{arsinh}x\)，它的反函数是双曲正弦函数
\[
\sinh x=\frac{e^x-e^{-x}}{2}.
\]
本题不是要求记忆函数名，而是要掌握有理化消根号的结构：
\[
\left(x+\sqrt{x^2+1}\right)\left(\sqrt{x^2+1}-x\right)=1.
\]

\textbf{可迁移方法：}

遇到
\[
\ln\left(x+\sqrt{x^2+a^2}\right)
\]
这类结构时，优先尝试对倒数有理化，构造 \(e^y\) 与 \(e^{-y}\)。若根式前后正负号互补，通常相加或相减可以消去根式或消去 \(x\)。

这类题的本质不是单纯“有理化技巧”，而是“制造第二个方程，再联立消元”的方程组思维。只不过第二个方程不是题目直接给的，而是由倒数有理化主动构造出来的。

\textbf{母题模板：}

\[
y=\ln\left(x+\sqrt{x^2+1}\right)
\quad\Longrightarrow\quad
\begin{cases}
e^y=x+\sqrt{x^2+1},\\
e^{-y}=\sqrt{x^2+1}-x,
\end{cases}
\]
所以
\[
x=\frac12(e^y-e^{-y}).
\]
\end{knowledgeBox}

\begin{mistakeBox}
\textbf{易错点：}

\begin{itemize}
  \item \textbf{把反函数误认为倒数：} \(f^{-1}(x)\) 是反函数，不是 \(\frac{1}{f(x)}\)。
  \item \textbf{硬平方导致计算膨胀：} 直接由 \(e^y=x+\sqrt{x^2+1}\) 移项平方也能做，但更容易漏项；本题最快入口是倒数有理化。
  \item \textbf{忘写定义域：} 反函数的定义域等于原函数的值域，本题为 \(\mathbb{R}\)。
  \item \textbf{对数合法性漏检：} 需要知道 \(x+\sqrt{x^2+1}>0\)，原函数才在全体实数上有定义。
\end{itemize}

\textbf{复盘提醒：}

看到 \(\ln(x+\sqrt{x^2+1})\)，要立即联想到“倒数有理化”和“双曲正弦”。考场上不要先平方硬解，先写出 \(e^y\) 与 \(e^{-y}\) 这对式子。
\end{mistakeBox}

\studySubsection{calc-01-composite-function}{复合函数的定义、定义域与拆解方法}
\knowledgeAnchor[MATH1-KN-CALC-0012]{复合函数}
\legacyKnowledgeAnchor{composite-function}
\noindent{\footnotesize\textbf{索引：} \knowledgeIndexRef[MATH1-KN-CALC-0012]{复合函数}}\par

\begin{knowledgeBox}
\textbf{概念定位：}

复合函数属于函数基础中的“输入输出链条”问题。它真正考的不是把 \(f(g(x))\) 这个符号背下来，而是看清楚：内层函数先把 \(x\) 变成中间变量 \(u\)，外层函数再把 \(u\) 变成最终函数值 \(y\)。

\textbf{定义：}

设函数 \(y=f(u)\) 的定义域为 \(D_1\)，函数 \(u=g(x)\) 在集合 \(D\) 上有定义，并且
\[
g(D)\subset D_1.
\]
则由
\[
y=f[g(x)],\qquad x\in D
\]
确定的函数称为由内层函数 \(u=g(x)\) 和外层函数 \(y=f(u)\) 构成的复合函数，其中 \(u\) 称为中间变量。

更一般地，如果没有预先给定集合 \(D\)，则复合函数 \(f\circ g\) 的自然定义域应写成
\[
D_{f\circ g}=\{x\in D_g\mid g(x)\in D_f\}.
\]
这句话非常重要：复合函数的定义域不是简单照抄内层函数定义域，而是“内层能算，且内层输出能送进外层”。

\textbf{直觉解释：}

复合函数就是一条流水线：
\[
x \xrightarrow{g} u=g(x) \xrightarrow{f} y=f(u).
\]
所以做复合函数题时，第一眼不是展开计算，而是先问两个问题：
\[
\text{第一，内层输出是什么？}
\qquad
\text{第二，这个输出落入外层函数的哪一段或哪个定义域？}
\]

\textbf{使用条件：}

\begin{itemize}
  \item 内层函数 \(g(x)\) 必须在当前 \(x\) 上有定义。
  \item 内层函数值 \(g(x)\) 必须属于外层函数 \(f\) 的定义域。
  \item 若外层函数是分段函数，判断分段条件时看的是内层输出 \(g(x)\)，不是原来的 \(x\)。
  \item 若由 \(f[\varphi(x)]\) 反求 \(\varphi(x)\)，要把 \(\varphi(x)\) 当成整体变量，并用附加条件确定分支。
\end{itemize}

\textbf{常见题型：}

\begin{itemize}
  \item 已知 \(f(x)\) 和 \(g(x)\)，求 \(f[g(x)]\)、\(g[f(x)]\)。
  \item 已知 \(f[\varphi(x)]\)，反求中间函数 \(\varphi(x)\)。
  \item 外层或内层为分段函数时，求复合函数表达式。
  \item 求复合函数的定义域和值域。
\end{itemize}

\textbf{常见误区：}

\begin{itemize}
  \item 只看原来的 \(x\) 属于哪一段，而忘记外层函数看的是内层输出。
  \item 已知 \(f[\varphi(x)]\) 时，把 \(x\) 直接代入 \(f(x)\)，没有把 \(\varphi(x)\) 当作整体。
  \item 求定义域时只管根号、分母和对数，却漏掉外层函数对内层输出的限制。
  \item 从平方关系反求中间函数时忘记正负分支。
\end{itemize}
\end{knowledgeBox}

\problemAnchor{MATH1-CALC-0004}
\begin{problemBox}
\begin{problemMeta}
\textbf{题目编号：} \problemRef{MATH1-CALC-0004} \quad
\textbf{索引：} \problemIndexRef{MATH1-CALC-0004} \quad
\textbf{来源：} 用户图片 / 教材例 1.4

\textbf{题型：} 已知 \(f[\varphi(x)]\) 反求中间函数 \quad
\textbf{难度：} 基础

\textbf{知识点：} \knowledgeRef[MATH1-KN-CALC-0012]{复合函数}，整体视角，定义域，值域，平方根分支
\end{problemMeta}

\textbf{题目：}

设
\[
f(x)=x^2,\qquad f[\varphi(x)]=-x^2+2x+3,
\]
且 \(\varphi(x)\ge 0\)，求 \(\varphi(x)\) 及其定义域和值域。
\end{problemBox}

\begin{solutionBox}
\textbf{思路入口：}

本题的关键是把 \(\varphi(x)\) 看成一个整体。因为外层函数是
\[
f(t)=t^2,
\]
所以
\[
f[\varphi(x)]=[\varphi(x)]^2.
\]
题目已经给出 \(f[\varphi(x)]\)，于是先得到关于整体 \(\varphi(x)\) 的平方方程，再由 \(\varphi(x)\ge 0\) 选择非负根。

\textbf{详细解答：}

由 \(f(x)=x^2\)，可知对任意中间变量 \(t\)，
\[
f(t)=t^2.
\]
令 \(t=\varphi(x)\)，则
\[
f[\varphi(x)]=[\varphi(x)]^2.
\]
由题设
\[
[\varphi(x)]^2=-x^2+2x+3.
\]
又因为题目给出 \(\varphi(x)\ge 0\)，所以只能取非负平方根：
\[
\varphi(x)=\sqrt{-x^2+2x+3}.
\]

\textbf{定义域：}

为了使 \(\varphi(x)\) 为实值函数，根号内必须非负：
\[
-x^2+2x+3\ge 0.
\]
等价于
\[
x^2-2x-3\le 0,
\]
即
\[
(x-3)(x+1)\le 0.
\]
所以
\[
-1\le x\le 3.
\]
故 \(\varphi(x)\) 的定义域为
\[
[-1,3].
\]

\textbf{值域：}

将根号内配方：
\[
-x^2+2x+3=-(x-1)^2+4.
\]
当 \(x\in[-1,3]\) 时，二次函数
\[
-(x-1)^2+4
\]
在 \(x=1\) 处取得最大值 \(4\)，在端点 \(x=-1,3\) 处取得最小值 \(0\)。因此
\[
-x^2+2x+3\in[0,4].
\]
开平方后
\[
\varphi(x)=\sqrt{-x^2+2x+3}\in[0,2].
\]
故 \(\varphi(x)\) 的值域为
\[
[0,2].
\]

\textbf{合法性检查：}

第一，\(\varphi(x)\ge 0\) 是选取正平方根的依据；如果没有这个条件，仅由 \([\varphi(x)]^2=-x^2+2x+3\) 不能唯一确定 \(\varphi(x)\)。第二，定义域必须来自根号内非负，而不是原外层函数 \(f(x)=x^2\) 的定义域，因为真正要求的是 \(\varphi(x)\) 这个函数的表达式。

\textbf{考场满分写法：}

必须写：
\[
f[\varphi(x)]=[\varphi(x)]^2=-x^2+2x+3.
\]
由 \(\varphi(x)\ge 0\)，得
\[
\varphi(x)=\sqrt{-x^2+2x+3}.
\]
再由
\[
-x^2+2x+3\ge 0
\]
得定义域
\[
[-1,3].
\]
最后配方
\[
-x^2+2x+3=-(x-1)^2+4
\]
得值域
\[
[0,2].
\]

\textbf{答案：}

\[
\boxed{\varphi(x)=\sqrt{-x^2+2x+3},\qquad D_\varphi=[-1,3],\qquad R_\varphi=[0,2].}
\]
\end{solutionBox}

\begin{knowledgeBox}
\textbf{核心知识点：}

已知 \(f[\varphi(x)]\) 且外层函数 \(f\) 明确时，要把 \(\varphi(x)\) 当成整体。若
\[
f(t)=t^2,
\]
则
\[
f[\varphi(x)]=[\varphi(x)]^2.
\]

\textbf{可迁移方法：}

遇到 \(f[\varphi(x)]\) 反求 \(\varphi(x)\)，按三步走：
\[
\text{识别外层函数}\ \longrightarrow\ \text{把内层整体代入}\ \longrightarrow\ \text{按附加条件选分支}.
\]
如果外层函数不是一一对应函数，例如平方函数，就必须格外关注分支条件。
\end{knowledgeBox}

\begin{mistakeBox}
\textbf{易错点：}

\begin{itemize}
  \item \textbf{漏掉正负分支：} 从 \([\varphi(x)]^2=G(x)\) 不能直接推出 \(\varphi(x)=\sqrt{G(x)}\)，本题能这样写是因为 \(\varphi(x)\ge 0\)。
  \item \textbf{定义域判断不完整：} \(\varphi(x)\) 的表达式含根号，必须要求 \(-x^2+2x+3\ge 0\)。
  \item \textbf{值域只看端点：} 根号内是开口向下的二次函数，最大值在 \(x=1\)，不能只代端点。
\end{itemize}

\textbf{复盘提醒：}

这题的母题是“已知外层函数和复合结果，反求内层函数”。核心动作不是展开，而是把内层函数当成一个整体未知量。
\end{mistakeBox}

\problemAnchor{MATH1-CALC-0005}
\begin{problemBox}
\begin{problemMeta}
\textbf{题目编号：} \problemRef{MATH1-CALC-0005} \quad
\textbf{索引：} \problemIndexRef{MATH1-CALC-0005} \quad
\textbf{来源：} 用户图片 / 教材例 1.5

\textbf{题型：} 分段函数复合 \quad
\textbf{难度：} 基础

\textbf{知识点：} \knowledgeRef[MATH1-KN-CALC-0012]{复合函数}，分段函数，内层函数取值范围，数形结合
\end{problemMeta}

\textbf{题目：}

设
\[
g(x)=
\begin{cases}
2-x, & x\le 0,\\
2+x, & x>0,
\end{cases}
\qquad
f(x)=
\begin{cases}
x^2, & x<0,\\
-x-1, & x\ge 0,
\end{cases}
\]
则求 \(g[f(x)]\)。
\end{problemBox}

\begin{solutionBox}
\textbf{思路入口：}

这是分段函数复合题。关键不是看原来的 \(x\) 是否大于 \(0\) 后直接套 \(g\)，而是先算内层函数
\[
u=f(x),
\]
再判断 \(u\) 落在 \(g\) 的哪一段。也就是说，外层 \(g\) 的分段条件看的是 \(f(x)\le 0\) 还是 \(f(x)>0\)，不是看原来的 \(x\le 0\) 还是 \(x>0\)。

\textbf{详细解答：}

外层函数 \(g\) 的分段规则是
\[
g(u)=
\begin{cases}
2-u, & u\le 0,\\
2+u, & u>0.
\end{cases}
\]
因此
\[
g[f(x)]=
\begin{cases}
2-f(x), & f(x)\le 0,\\
2+f(x), & f(x)>0.
\end{cases}
\]
现在判断 \(f(x)\) 的符号。

当 \(x\ge 0\) 时，
\[
f(x)=-x-1\le -1<0.
\]
所以此时 \(f(x)\le 0\)，应代入 \(g\) 的第一段：
\[
g[f(x)]=2-f(x)=2-(-x-1)=x+3.
\]

当 \(x<0\) 时，
\[
f(x)=x^2>0.
\]
所以此时 \(f(x)>0\)，应代入 \(g\) 的第二段：
\[
g[f(x)]=2+f(x)=2+x^2.
\]
综上，
\[
g[f(x)]=
\begin{cases}
2+x^2, & x<0,\\
3+x, & x\ge 0.
\end{cases}
\]

\textbf{合法性检查：}

内层函数 \(f(x)\) 在全体实数上有定义，外层函数 \(g(x)\) 也在全体实数上有定义，所以复合函数 \(g[f(x)]\) 的定义域为
\[
\mathbb{R}.
\]
边界点 \(x=0\) 必须归入 \(f\) 的第二段，因为题设为 \(x\ge 0\)。此时
\[
f(0)=-1\le 0,
\]
所以外层 \(g\) 取第一段，得到
\[
g[f(0)]=3.
\]
这与公式 \(3+x\) 在 \(x=0\) 的取值一致。

\textbf{考场满分写法：}

必须先写：
\[
g[f(x)]=
\begin{cases}
2-f(x), & f(x)\le 0,\\
2+f(x), & f(x)>0.
\end{cases}
\]
再判断：
\[
f(x)\le 0\Longleftrightarrow x\ge 0,\qquad
f(x)>0\Longleftrightarrow x<0.
\]
最后代回得
\[
g[f(x)]=
\begin{cases}
2+x^2, & x<0,\\
3+x, & x\ge 0.
\end{cases}
\]

\textbf{答案：}

\[
\boxed{
g[f(x)]=
\begin{cases}
2+x^2, & x<0,\\
3+x, & x\ge 0.
\end{cases}}
\]
\end{solutionBox}

\begin{knowledgeBox}
\textbf{核心知识点：}

分段复合函数的核心是“内层决定外层落段”。若要求 \(g[f(x)]\)，则 \(g\) 的分段条件要用 \(f(x)\) 去判断。

\textbf{可迁移方法：}

遇到分段函数复合，按四步走：
\[
\text{写外层分段规则}\ \longrightarrow\ \text{把外层输入换成内层函数}\ \longrightarrow\ \text{判断内层函数落段}\ \longrightarrow\ \text{代回化简}.
\]
如果内层函数本身也是分段函数，建议先画图或列符号表。图像能快速告诉我们内层输出落在正半轴、负半轴还是跨越分段点。
\end{knowledgeBox}

\begin{mistakeBox}
\textbf{易错点：}

\begin{itemize}
  \item \textbf{看错分段变量：} 求 \(g[f(x)]\) 时，\(g\) 的分段依据是 \(f(x)\)，不是原来的 \(x\)。
  \item \textbf{边界点归属错误：} \(x=0\) 应先归入 \(f(x)=-x-1\)，再判断 \(f(0)=-1\le 0\)。
  \item \textbf{分段顺序混乱：} 先判断内层 \(f(x)\) 的符号，再代入外层 \(g\)，不要把两层函数的分段条件揉在一起。
\end{itemize}

\textbf{复盘提醒：}

这题要练成条件反射：复合函数先看内层输出。外层函数问“我的输入落在哪一段”，这个输入不是 \(x\)，而是 \(f(x)\)。
\end{mistakeBox}

\studySubsection{calc-01-implicit-function}{隐函数的概念与观察法}
\knowledgeAnchor[MATH1-KN-CALC-0029]{隐函数}
\legacyKnowledgeAnchor{implicit-function}
\noindent{\footnotesize\textbf{索引：} \knowledgeIndexRef[MATH1-KN-CALC-0029]{隐函数}}\par

\begin{knowledgeBox}
\textbf{概念定位：}

隐函数属于“变量关系”问题。显函数直接写成 \(y=f(x)\)，输入 \(x\) 后能直接算出 \(y\)；隐函数则先给出一个方程
\[
F(x,y)=0,
\]
它把 \(x,y\) 捆在一起。若在某个区间内每个 \(x\) 都能唯一确定一个 \(y\)，则这个方程在该区间上确定了一个隐函数
\[
y=y(x).
\]

\textbf{定义：}

设方程 \(F(x,y)=0\)。如果当 \(x\) 在某区间 \(I\) 内取任一值时，总存在唯一的 \(y\) 满足该方程，则称方程 \(F(x,y)=0\) 在该区间内确定了一个隐函数 \(y=y(x)\)。

\textbf{为什么“唯一”是关键：}

隐函数不是说方程里同时有 \(x,y\) 就行，而是要由 \(x\) 唯一反推出 \(y\)。如果某个 \(x\) 对应两个不同的 \(y\)，那描述的是多值关系，不能直接称为一个函数 \(y=y(x)\)。这和反函数中的“一一对应”是同一种底层思想。

\textbf{显化与不易显化：}

有些隐函数容易显化。例如
\[
x+y^3-1=0
\]
可以改写为
\[
y=\sqrt[3]{1-x}.
\]
有些隐函数不容易显化，例如
\[
\sin(xy)=\ln\frac{x+e}{y}+1.
\]
这类方程很难整理成 \(y=y(x)\) 的显式表达，但在局部仍可能确定隐函数。后续学习隐函数求导时，常常不需要显化，只要直接对方程两边求导。

\textbf{特殊点求值的观察法：}

如果只要求某个点的函数值 \(y(x_0)\)，不一定要解出整个 \(y=y(x)\)。把 \(x=x_0\) 代入方程后，若关于 \(y\) 的方程容易观察出唯一解，就直接求该点函数值。

例如，方程
\[
\ln y-\frac{x}{y}+x=0,\qquad x>0
\]
确定 \(y=y(x)\)。当 \(x=2\) 时，方程化为
\[
\ln y-\frac{2}{y}+2=0.
\]
令
\[
h(y)=\ln y-\frac{2}{y}+2,\qquad y>0.
\]
直接观察 \(y=1\) 时
\[
h(1)=0-2+2=0.
\]
若要补足唯一性，可看
\[
h'(y)=\frac1y+\frac{2}{y^2}>0,
\]
所以 \(h(y)\) 在 \(y>0\) 上严格递增，零点唯一。因此
\[
y(2)=1.
\]

再如，方程
\[
\ln y+e^{y-1}=\frac{x}{2}
\]
确定 \(y=y(x)\)。当 \(x=2\) 时，
\[
\ln y+e^{y-1}=1.
\]
令
\[
k(y)=\ln y+e^{y-1},\qquad y>0.
\]
观察得
\[
k(1)=0+1=1.
\]
又因为
\[
k'(y)=\frac1y+e^{y-1}>0,
\]
所以解唯一，故
\[
y(2)=1.
\]

\textbf{常见误区：}

\begin{itemize}
  \item 看到 \(F(x,y)=0\) 就默认一定有隐函数，忽略了唯一性。
  \item 为了求一个点 \(y(x_0)\)，强行解出整条函数表达式，计算反而变繁。
  \item 混淆“隐函数存在”和“能显式写出”。很多隐函数可以存在，但未必容易显化。
\end{itemize}
\end{knowledgeBox}

\studySubsection{calc-01-boundedness}{函数的有界性}
\knowledgeAnchor[MATH1-KN-CALC-0017]{有界性}
\legacyKnowledgeAnchor{boundedness}
\noindent{\footnotesize\textbf{索引：} \knowledgeIndexRef[MATH1-KN-CALC-0017]{有界性}}\par

\begin{knowledgeBox}
\textbf{概念定位：}

有界性回答的是：函数值在某个指定区间内会不会无限跑远。它不是单个点的性质，而是函数在一个集合或区间上的整体性质。

\textbf{定义：}

设 \(f(x)\) 的定义域为 \(D\)，数集 \(I\subset D\)。如果存在正数 \(M\)，使得对任意 \(x\in I\)，都有
\[
|f(x)|\le M,
\]
则称 \(f(x)\) 在 \(I\) 上有界；如果这样的 \(M\) 不存在，则称 \(f(x)\) 在 \(I\) 上无界。

\textbf{直觉解释：}

几何上，若在区间 \(I\) 内，函数图像能被两条水平直线
\[
y=M,\qquad y=-M
\]
夹住，则函数在 \(I\) 上有界。

\textbf{使用条件：}

讨论有界性必须先明确区间 \(I\)。同一个函数在不同区间上有界性可能不同，例如
\[
y=\frac1x
\]
在 \((2,+\infty)\) 上有界，但在 \((0,2)\) 上无界。

\textbf{无界的常见判据：}

如果在区间 \(I\) 内或端点附近存在点 \(x_0\)，使得
\[
\lim_{x\to x_0}f(x)=\infty
\quad\text{或}\quad
\lim_{x\to x_0}f(x)=-\infty,
\]
则函数在该区间上无界。因为任意大的水平线都无法把图像夹住。
\end{knowledgeBox}

\problemAnchor{MATH1-CALC-0006}
\begin{problemBox}
\begin{problemMeta}
\textbf{题目编号：} \problemRef{MATH1-CALC-0006} \quad
\textbf{索引：} \problemIndexRef{MATH1-CALC-0006} \quad
\textbf{来源：} 用户图片 / 教材例 1.6

\textbf{题型：} 证明函数有界 \quad
\textbf{难度：} 基础

\textbf{知识点：} \knowledgeRef[MATH1-KN-CALC-0017]{有界性}，绝对值，基本不等式，放缩
\end{problemMeta}

\textbf{题目：}

证明函数
\[
f(x)=\frac{x}{1+x^2}
\]
在 \((-\infty,+\infty)\) 内有界。
\end{problemBox}

\begin{solutionBox}
\textbf{思路入口：}

证明有界的目标是找到一个常数 \(M\)，使得
\[
|f(x)|\le M
\]
对所有实数 \(x\) 都成立。本题分母 \(1+x^2\) 永远为正，关键是把
\[
\left|\frac{x}{1+x^2}\right|
\]
放缩到一个常数。自然想到基本不等式或
\[
1+x^2\ge 2|x|.
\]

\textbf{详细解答一：用基本不等式}

当 \(x=0\) 时，
\[
f(0)=0.
\]
当 \(x\ne 0\) 时，
\[
|f(x)|=\frac{|x|}{1+x^2}.
\]
分子分母同除以 \(|x|\)，得
\[
|f(x)|=\frac{1}{\frac1{|x|}+|x|}.
\]
由基本不等式
\[
a+\frac1a\ge 2,\qquad a>0,
\]
令 \(a=|x|\)，得
\[
\frac1{|x|}+|x|\ge 2.
\]
因此
\[
|f(x)|\le \frac12.
\]
结合 \(x=0\) 的情况，对任意 \(x\in\mathbb{R}\)，都有
\[
|f(x)|\le \frac12.
\]
所以 \(f(x)\) 在 \((-\infty,+\infty)\) 内有界。

\textbf{详细解答二：直接放缩}

由
\[
\frac{1+x^2}{2}\ge \sqrt{x^2}=|x|
\]
可得
\[
1+x^2\ge 2|x|.
\]
当 \(x\ne0\) 时，
\[
|f(x)|=\frac{|x|}{1+x^2}\le \frac{|x|}{2|x|}=\frac12.
\]
当 \(x=0\) 时，仍有 \(f(0)=0\)。综上，
\[
|f(x)|\le\frac12,\qquad x\in\mathbb{R}.
\]

\textbf{合法性检查：}

第一，分母 \(1+x^2>0\)，所以函数在全体实数上有定义。第二，第一种方法中除以 \(|x|\) 前必须单独处理 \(x=0\)。第三，证明有界不要求找到最小的 \(M\)，只要找到一个可行的上界即可；本题实际可取 \(M=\frac12\)。

\textbf{考场满分写法：}

写出
\[
1+x^2\ge 2|x|.
\]
当 \(x\ne0\) 时，
\[
|f(x)|=\frac{|x|}{1+x^2}\le \frac{|x|}{2|x|}=\frac12.
\]
当 \(x=0\) 时，\(f(0)=0\)。故对任意 \(x\in\mathbb{R}\)，有
\[
|f(x)|\le \frac12,
\]
所以 \(f(x)\) 在 \((-\infty,+\infty)\) 内有界。

\textbf{答案：}

\[
\boxed{|f(x)|\le \frac12,\quad x\in\mathbb{R},\ \text{故 } f(x)\text{ 有界}.}
\]
\end{solutionBox}

\begin{knowledgeBox}
\textbf{核心知识点：}

证明有界的本质是找 \(M\)。对于有理函数
\[
\frac{x}{1+x^2},
\]
常用入口是
\[
1+x^2\ge 2|x|.
\]

\textbf{可迁移方法：}

遇到形如
\[
\frac{|x|}{1+x^2}
\]
的结构，优先考虑把分母放缩为 \(2|x|\)。如果分母分子都含 \(x\)，还要特别注意 \(x=0\) 是否需要单独处理。

\textbf{为什么放缩优先于极值法：}

本题当然也可以用后续导数方法：先讨论
\[
f(x)=\frac{x}{1+x^2}
\]
的单调性，再找极值点，最后得到最大绝对值。但在“只需证明有界”的题里，这条路过重。放缩法直接给出
\[
|f(x)|\le \frac12,
\]
等于绕开了“求导、分区间、找极值、比较端点”的整套流程。考场上，若题目只问“有界”或“估值”，优先想放缩；若题目问“最大值、最小值、取等条件”，再考虑极值法。
\end{knowledgeBox}

\begin{mistakeBox}
\textbf{易错点：}

\begin{itemize}
  \item \textbf{忘记先取绝对值：} 有界性通常用 \(|f(x)|\le M\) 表示，而不是只证明 \(f(x)\le M\)。
  \item \textbf{除以 \(|x|\) 时漏掉 \(x=0\)：} \(|x|=0\) 时不能作为除数。
  \item \textbf{误以为必须求最小上界：} 证明有界只需找到一个 \(M\)，不一定要求最小。
\end{itemize}

\textbf{复盘提醒：}

看到“证明有界”，先把目标翻译成“找 \(M\)”。本题要记住标准放缩：
\[
1+x^2\ge 2|x|.
\]
\end{mistakeBox}

\studySubsection{calc-01-boundedness-practice}{有界性放缩练习组}

\problemAnchor{MATH1-CALC-0010}
\begin{problemBox}
\begin{problemMeta}
\textbf{题目编号：} \problemRef{MATH1-CALC-0010} \quad
\textbf{索引：} \problemIndexRef{MATH1-CALC-0010} \quad
\textbf{来源：} 用户要求生成 / 放缩练习

\textbf{题型：} 放缩证明有界与估值 \quad
\textbf{难度：} 基础到中等

\textbf{知识点：} \knowledgeRef[MATH1-KN-CALC-0017]{有界性}，基本不等式，分母放缩，有理化
\end{problemMeta}

\textbf{练习要求：}

以下题目尽量不要先求导。先把目标翻译成“找一个统一上界”，再尝试用放缩完成。

\begin{enumerate}
  \item 证明函数
  \[
  f(x)=\frac{x^2}{1+x^4}
  \]
  在 \((-\infty,+\infty)\) 内有界。

  \item 证明函数
  \[
  f(x)=\frac{|x|}{1+x^2+x^4}
  \]
  在 \((-\infty,+\infty)\) 内有界。

  \item 证明对任意 \(x\in\mathbb{R}\)，有
  \[
  0<\sqrt{x^2+1}-|x|\le 1.
  \]

  \item 证明函数
  \[
  f(x)=\frac{x}{\sqrt{x^2+1}+1}
  \]
  在 \((-\infty,+\infty)\) 内有界。

  \item 设 \(x>0\)，证明
  \[
  \frac{x}{(1+x)^2}\le \frac14.
  \]
\end{enumerate}
\end{problemBox}

\begin{solutionBox}
\textbf{提示与答案：}

\begin{enumerate}
  \item 用
  \[
  1+x^4\ge 2x^2
  \]
  得
  \[
  0\le \frac{x^2}{1+x^4}\le \frac12.
  \]

  \item 分母越大，正分式越小。可先丢掉 \(x^4\)，用
  \[
  1+x^2+x^4\ge 1+x^2\ge 2|x|
  \]
  得
  \[
  0\le \frac{|x|}{1+x^2+x^4}\le \frac12.
  \]
  注意 \(x=0\) 时原式为 \(0\)，不需要除以 \(|x|\)。

  \item 用有理化：
  \[
  \sqrt{x^2+1}-|x|
  =\frac{1}{\sqrt{x^2+1}+|x|}.
  \]
  分母大于等于 \(1\)，所以
  \[
  0<\sqrt{x^2+1}-|x|\le 1.
  \]

  补充：若在分情况时遇到
  \[
  \sqrt{x^2+1}+x<1,
  \]
  必须先说明这是在 \(x<0\) 的情况下讨论。因为当 \(x<0\) 时，\(|x|=-x\)，所以
  \[
  \sqrt{x^2+1}+x=\sqrt{x^2+1}-|x|.
  \]
  用有理化可得
  \[
  \sqrt{x^2+1}+x
  =\frac{1}{\sqrt{x^2+1}-x}.
  \]
  此时 \(x<0\)，故 \(-x>0\)，分母
  \[
  \sqrt{x^2+1}-x=\sqrt{x^2+1}+|x|>1,
  \]
  所以
  \[
  0<\sqrt{x^2+1}+x<1.
  \]
  注意：若不加 \(x<0\) 这个条件，原不等式不成立；例如 \(x=0\) 时左边等于 \(1\)。

  \item 取绝对值：
  \[
  |f(x)|=\frac{|x|}{\sqrt{x^2+1}+1}.
  \]
  因为 \(\sqrt{x^2+1}>|x|\)，所以
  \[
  |f(x)|<1.
  \]
  因此 \(f(x)\) 有界。

  \item 用
  \[
  (1+x)^2\ge 4x,\qquad x>0,
  \]
  得
  \[
  \frac{x}{(1+x)^2}\le \frac14.
  \]
\end{enumerate}

\textbf{复盘重点：}

第 1、2、5 题练“分母放大，正分式变小”；第 3 题练“根式差有理化”；第 4 题练“根式分母天然大于 \(|x|\)”。这些都是有界性题中替代求导极值法的常用入口。
\end{solutionBox}

\begin{mistakeBox}
\textbf{易错点：}

\begin{itemize}
  \item 正分式中，分母放大后分式变小；如果分母或分子可能为负，必须先处理符号。
  \item 含 \(|x|\) 的放缩常要单独检查 \(x=0\)，避免非法约分。
  \item 根式差 \(\sqrt{x^2+1}-|x|\) 不要硬估，优先有理化。
\end{itemize}

\textbf{复盘提醒：}

做完后不要只看答案对不对，要给每题标出“触发信号”：分母有 \(1+x^4\)、分母有 \(1+x^2\)、根式差、根式分母、\((1+x)^2\)。以后看到这些结构，直接调用对应放缩模板。
\end{mistakeBox}

\studySubsection{calc-01-monotonicity}{函数的单调性}
\knowledgeAnchor[MATH1-KN-CALC-0021]{单调性}
\legacyKnowledgeAnchor{monotonicity}
\noindent{\footnotesize\textbf{索引：} \knowledgeIndexRef[MATH1-KN-CALC-0021]{单调性}}\par

\begin{knowledgeBox}
\textbf{概念定位：}

单调性描述函数在区间上随自变量增大时的整体变化趋势。它仍然是区间性质，不是单点性质。

\textbf{定义：}

设 \(f(x)\) 的定义域为 \(D\)，区间 \(I\subset D\)。如果对任意 \(x_1,x_2\in I\)，当 \(x_1<x_2\) 时恒有
\[
f(x_1)<f(x_2),
\]
则称 \(f(x)\) 在 \(I\) 上严格单调增加。

如果对任意 \(x_1,x_2\in I\)，当 \(x_1<x_2\) 时恒有
\[
f(x_1)>f(x_2),
\]
则称 \(f(x)\) 在 \(I\) 上严格单调减少。

\textbf{乘积判别式：}

对任意 \(x_1,x_2\in I\)，\(x_1\ne x_2\)，有
\[
f(x)\text{ 严格单调增加}
\quad\Longleftrightarrow\quad
(x_1-x_2)[f(x_1)-f(x_2)]>0.
\]
有
\[
f(x)\text{ 严格单调减少}
\quad\Longleftrightarrow\quad
(x_1-x_2)[f(x_1)-f(x_2)]<0.
\]

\textbf{为什么这个判别式成立：}

若 \(x_1<x_2\)，则 \(x_1-x_2<0\)。严格递增时 \(f(x_1)<f(x_2)\)，所以
\[
f(x_1)-f(x_2)<0.
\]
两个负数相乘为正，故乘积大于 \(0\)。这正是“自变量差”和“函数值差”同号的意思。

\textbf{常见误区：}

\begin{itemize}
  \item 单调性必须说明在哪个区间上讨论。
  \item 后续常用导数判断单调性，但定义法和乘积判别式在没有导数时同样重要。
  \item 图像变换会改变单调性：\(f(-x)\) 通常与 \(f(x)\) 单调性相反，而 \(-f(-x)\) 又会再反一次。
\end{itemize}
\end{knowledgeBox}

\problemAnchor{MATH1-CALC-0007}
\begin{problemBox}
\begin{problemMeta}
\textbf{题目编号：} \problemRef{MATH1-CALC-0007} \quad
\textbf{索引：} \problemIndexRef{MATH1-CALC-0007} \quad
\textbf{来源：} 用户图片 / 教材例 1.7

\textbf{题型：} 单调性与图像变换 \quad
\textbf{难度：} 中等

\textbf{知识点：} \knowledgeRef[MATH1-KN-CALC-0021]{单调性}定义，图像对称，复合与变号
\end{problemMeta}

\textbf{题目：}

设 \(f(x)\) 在 \((-\infty,+\infty)\) 上有定义，任给 \(x_1,x_2\)，\(x_1\ne x_2\)，均有
\[
(x_1-x_2)\bigl[f(x_1)-f(x_2)\bigr]>0.
\]
则以下函数一定单调增加的是
\[
\text{(A) } |f(x)|,\qquad
\text{(B) } f(|x|),\qquad
\text{(C) } f(-x),\qquad
\text{(D) } -f(-x).
\]
\end{problemBox}

\begin{solutionBox}
\textbf{思路入口：}

题设给出的不是导数，而是单调性的乘积判别式。它表示 \(f(x)\) 在全体实数上严格单调增加。接下来要看四个变换对单调性的影响：取绝对值、输入取绝对值、横向反射、先横向反射再纵向反射。

\textbf{详细解答：}

由题设
\[
(x_1-x_2)\bigl[f(x_1)-f(x_2)\bigr]>0
\]
可知 \(f(x)\) 是严格单调增加函数。

先看选项 (D)。令
\[
F(x)=-f(-x).
\]
任取 \(x_1<x_2\)，则
\[
-x_1>-x_2.
\]
因为 \(f\) 严格单调增加，所以
\[
f(-x_1)>f(-x_2).
\]
两边同时乘以 \(-1\)，不等号反向：
\[
-f(-x_1)<-f(-x_2).
\]
即
\[
F(x_1)<F(x_2).
\]
因此
\[
-f(-x)
\]
一定严格单调增加。

\textbf{为什么其他选项不一定：}

选项 (C) 中，\(f(-x)\) 是把 \(f(x)\) 的图像关于 \(y\) 轴对称。横向反射会把严格递增变成严格递减，所以不符合。

选项 (A) 中，\(|f(x)|\) 会把 \(f(x)<0\) 的部分翻到 \(x\) 轴上方，不一定保持单调。例如取 \(f(x)=x\)，则
\[
|f(x)|=|x|
\]
在 \(\mathbb{R}\) 上不是单调增加。

选项 (B) 中，\(f(|x|)\) 会先把 \(x<0\) 的部分折到 \(x>0\)，通常得到偶函数，不可能在整个实数轴上严格单调增加。例如取 \(f(x)=x\)，则
\[
f(|x|)=|x|
\]
仍不是单调增加。

\textbf{合法性检查：}

题设中 \(f\) 在全体实数上有定义，所以 \(-x\)、\(|x|\) 作为输入时都不会超出定义域。判断 (D) 时不能只凭图像直觉，最好用定义法写出 \(x_1<x_2\Rightarrow F(x_1)<F(x_2)\)。

\textbf{考场满分写法：}

由题设可知 \(f\) 严格单调增加。令
\[
F(x)=-f(-x).
\]
若 \(x_1<x_2\)，则 \(-x_1>-x_2\)。由 \(f\) 递增得
\[
f(-x_1)>f(-x_2).
\]
故
\[
-f(-x_1)<-f(-x_2),
\]
即 \(F(x_1)<F(x_2)\)，所以 \(F(x)=-f(-x)\) 严格单调增加。

\textbf{答案：}

\[
\boxed{\text{D}}
\]
\end{solutionBox}

\begin{knowledgeBox}
\textbf{核心知识点：}

若 \(f(x)\) 严格递增，则
\[
f(-x)
\]
严格递减，而
\[
-f(-x)
\]
严格递增。图像上看，\(f(-x)\) 是关于 \(y\) 轴对称，\(-f(-x)\) 是关于原点对称。

\textbf{可迁移方法：}

遇到单调性与图像变换题，优先用定义法验证。设新函数为 \(F(x)\)，任取 \(x_1<x_2\)，追踪输入变换后不等号如何变化，再判断 \(F(x_1)\) 与 \(F(x_2)\) 的大小。
\end{knowledgeBox}

\begin{mistakeBox}
\textbf{易错点：}

\begin{itemize}
  \item \textbf{把 \(|f(x)|\) 误认为保持单调：} 取绝对值会翻折负值部分，可能破坏单调性。
  \item \textbf{把 \(f(|x|)\) 误认为只限制定义域：} 输入取绝对值会把左半轴折到右半轴，整体通常不是单调函数。
  \item \textbf{漏掉两次反向：} \(x\mapsto -x\) 反一次单调性，外面再乘 \(-1\) 又反一次，所以 \(-f(-x)\) 保持递增。
\end{itemize}

\textbf{复盘提醒：}

单调性题不要只看形状，要能用定义法写出不等号链条。尤其是 \(f(-x)\)、\(-f(x)\)、\(-f(-x)\) 三类变换，要分别记清。
\end{mistakeBox}

\studySubsection{calc-01-parity}{函数的奇偶性}
\knowledgeAnchor[MATH1-KN-CALC-0024]{奇函数}
\legacyKnowledgeAnchor{odd-function}
\noindent{\footnotesize\textbf{索引：} \knowledgeIndexRef[MATH1-KN-CALC-0024]{奇函数}}\par

\begin{knowledgeBox}
\textbf{概念定位：}

奇偶性是函数图像对称性的代数表达。偶函数图像关于 \(y\) 轴对称，奇函数图像关于原点对称。

\textbf{定义：}

设函数 \(f(x)\) 的定义域 \(D\) 关于原点对称，即
\[
x\in D\quad\Longrightarrow\quad -x\in D.
\]
如果对任意 \(x\in D\)，恒有
\[
f(-x)=f(x),
\]
则称 \(f(x)\) 为偶函数。

如果对任意 \(x\in D\)，恒有
\[
f(-x)=-f(x),
\]
则称 \(f(x)\) 为奇函数。

\textbf{判断奇偶性的第一步：}

必须先检查定义域是否关于原点对称。若定义域不关于原点对称，通常不能谈奇偶性。

\textbf{常用结构：}

对任意函数 \(f(x)\)，在定义域允许的前提下，
\[
f(x)+f(-x)
\]
是偶函数，
\[
f(x)-f(-x)
\]
是奇函数。

进一步令
\[
u(x)=\frac12[f(x)+f(-x)],\qquad
v(x)=\frac12[f(x)-f(-x)],
\]
则 \(u(x)\) 是偶函数，\(v(x)\) 是奇函数，且
\[
f(x)=u(x)+v(x).
\]
这说明任意函数都可以分解成“偶函数部分 + 奇函数部分”。

\textbf{复合函数的奇偶性规则：}

记忆口诀：
\[
\text{内层为偶，整体为偶；内层为奇，整体看外层。}
\]
也就是说，若 \(g(x)\) 是偶函数，则 \(f[g(x)]\) 通常为偶函数；若 \(g(x)\) 是奇函数，则外层 \(f\) 为偶时整体为偶，外层 \(f\) 为奇时整体为奇。

\textbf{导数与积分中的奇偶互换：}

若 \(f(x)\) 可导，则求导一次会使奇偶性互换：
\[
f\text{ 奇}\Rightarrow f'\text{ 偶},\qquad
f\text{ 偶}\Rightarrow f'\text{ 奇}.
\]
若
\[
F(x)=\int_0^x f(t)\,dt,
\]
则
\[
f\text{ 奇}\Rightarrow F\text{ 偶},\qquad
f\text{ 偶}\Rightarrow F\text{ 奇}.
\]
\end{knowledgeBox}

\problemAnchor{MATH1-CALC-0008}
\begin{problemBox}
\begin{problemMeta}
\textbf{题目编号：} \problemRef{MATH1-CALC-0008} \quad
\textbf{索引：} \problemIndexRef{MATH1-CALC-0008} \quad
\textbf{来源：} 用户图片 / 教材例 1.8

\textbf{题型：} 用定义证明奇函数 \quad
\textbf{难度：} 基础

\textbf{知识点：} \knowledgeRef[MATH1-KN-CALC-0024]{奇函数}，函数方程，赋值法
\end{problemMeta}

\textbf{题目：}

设对任意 \(x,y\)，都有
\[
f(x+y)=f(x)+f(y),
\]
证明：\(f(x)\) 是奇函数。
\end{problemBox}

\begin{solutionBox}
\textbf{思路入口：}

要证明 \(f(x)\) 是奇函数，目标是证明
\[
f(-x)=-f(x).
\]
题设是对任意 \(x,y\) 成立的函数方程，所以自然使用赋值法：先令 \(x=y=0\) 求 \(f(0)\)，再令 \(y=-x\) 建立 \(f(x)\) 与 \(f(-x)\) 的关系。

\textbf{详细解答：}

令 \(x=0,y=0\)，由题设
\[
f(0+0)=f(0)+f(0).
\]
即
\[
f(0)=2f(0),
\]
所以
\[
f(0)=0.
\]

再令 \(y=-x\)，则
\[
f(x+(-x))=f(x)+f(-x).
\]
左边为
\[
f(0)=0,
\]
因此
\[
0=f(x)+f(-x).
\]
所以
\[
f(-x)=-f(x).
\]
这正是奇函数的定义，故 \(f(x)\) 是奇函数。

\textbf{合法性检查：}

题目说“任意 \(x,y\)”都成立，说明函数定义域至少对这些加法操作封闭，赋值 \(x=y=0\) 与 \(y=-x\) 是合法的。证明奇函数时还需要默认定义域关于原点对称；本题“任意 \(x,y\)”的条件已保证 \(-x\) 可以作为自变量参与讨论。

\textbf{考场满分写法：}

令 \(x=y=0\)，得
\[
f(0)=2f(0),
\]
故
\[
f(0)=0.
\]
再令 \(y=-x\)，得
\[
f(0)=f(x)+f(-x).
\]
于是
\[
f(-x)=-f(x),
\]
故 \(f(x)\) 为奇函数。

\textbf{答案：}

\[
\boxed{f(-x)=-f(x),\ \text{故 } f(x)\text{ 是奇函数}.}
\]
\end{solutionBox}

\begin{knowledgeBox}
\textbf{核心知识点：}

函数方程中出现“任意 \(x,y\)”时，优先尝试赋特殊值。证明奇函数常用
\[
x+y=0
\]
来制造 \(f(0)\)。

\textbf{可迁移方法：}

遇到
\[
f(x+y)=f(x)+f(y)
\]
这类加法型函数方程，先求 \(f(0)\)，再用 \(y=-x\) 得到奇偶性或相反数关系。
\end{knowledgeBox}

\begin{mistakeBox}
\textbf{易错点：}

\begin{itemize}
  \item \textbf{直接令 \(y=-x\) 但忘记先求 \(f(0)\)：} 需要先知道 \(f(0)=0\) 才能得到 \(f(x)+f(-x)=0\)。
  \item \textbf{把奇函数误写成 \(f(-x)=f(x)\)：} 这是偶函数条件。
  \item \textbf{忽略“任意”条件的作用：} 正因为题设对任意 \(x,y\) 成立，才可以自由赋值。
\end{itemize}

\textbf{复盘提醒：}

函数方程题的第一反应是“赋值”。常用三件套：令变量为 \(0\)，令两个变量互为相反数，令两个变量相等。
\end{mistakeBox}

\studySubsection{calc-01-periodicity}{函数的周期性}
\knowledgeAnchor[MATH1-KN-CALC-0026]{周期函数}
\legacyKnowledgeAnchor{periodic-function}
\noindent{\footnotesize\textbf{索引：} \knowledgeIndexRef[MATH1-KN-CALC-0026]{周期函数}}\par

\begin{knowledgeBox}
\textbf{概念定位：}

周期性描述函数值按固定长度重复出现。它的本质是“平移不变”：
\[
f(x+T)=f(x).
\]

\textbf{定义：}

设函数 \(f(x)\) 的定义域为 \(D\)。如果存在非零常数 \(T\)，使得对任意 \(x\in D\)，都有 \(x+T\in D\)，且
\[
f(x+T)=f(x),
\]
则称 \(f(x)\) 为周期函数，\(T\) 称为 \(f(x)\) 的一个周期。通常所说“周期”多指最小正周期，但证明题中只要证明某个 \(T\ne0\) 是周期即可。

\textbf{重要结论：}

\begin{itemize}
  \item 若 \(f(x)\) 以 \(T\) 为周期，则 \(f(ax+b)\) 以 \(\dfrac{T}{|a|}\) 为周期，其中 \(a\ne0\)。
  \item 若 \(g(x)\) 是周期函数，则复合函数 \(f[g(x)]\) 也是周期函数。
  \item 若 \(f(x)\) 是以 \(T\) 为周期的可导函数，则 \(f'(x)\) 也以 \(T\) 为周期。
  \item 若 \(f(x)\) 是以 \(T\) 为周期的连续函数，则
  \[
  F(x)=\int_0^x f(t)\,dt
  \]
  以 \(T\) 为周期的充要条件是
  \[
  \int_0^T f(x)\,dx=0.
  \]
\end{itemize}

\textbf{常见误区：}

\begin{itemize}
  \item 证明 \(T\) 是周期，不等于证明它是最小正周期。
  \item 递推式中出现 \(x-\pi\) 或 \(x+\pi\) 时，要多次使用等式，把函数值推回到同一个位置。
  \item 看到 \(\sin(x+\pi)=-\sin x\)、\(\sin(x+2\pi)=\sin x\) 时，要主动联想到抵消结构。
\end{itemize}
\end{knowledgeBox}

\problemAnchor{MATH1-CALC-0009}
\begin{problemBox}
\begin{problemMeta}
\textbf{题目编号：} \problemRef{MATH1-CALC-0009} \quad
\textbf{索引：} \problemIndexRef{MATH1-CALC-0009} \quad
\textbf{来源：} 用户图片 / 教材例 1.9

\textbf{题型：} 由递推关系证明周期性 \quad
\textbf{难度：} 基础

\textbf{知识点：} \knowledgeRef[MATH1-KN-CALC-0026]{周期函数}，递推式，三角函数诱导公式
\end{problemMeta}

\textbf{题目：}

设函数 \(f(x)\) 在 \((-\infty,+\infty)\) 上满足
\[
f(x)=f(x-\pi)+\sin x.
\]
证明：\(f(x)\) 是以 \(T=2\pi\) 为周期的周期函数。
\end{problemBox}

\begin{solutionBox}
\textbf{思路入口：}

要证明 \(2\pi\) 是周期，只需证明
\[
f(x+2\pi)=f(x).
\]
题设给的是一步后退 \(\pi\) 的关系，所以对 \(x+2\pi\)、\(x+\pi\) 连续使用题设，把 \(f(x+2\pi)\) 逐步推回 \(f(x)\)。中间会出现
\[
\sin(x+\pi)=-\sin x,\qquad \sin(x+2\pi)=\sin x,
\]
正好相互抵消。

\textbf{详细解答：}

由题设，对任意实数 \(t\)，有
\[
f(t)=f(t-\pi)+\sin t.
\]
令 \(t=x+2\pi\)，得
\[
f(x+2\pi)=f(x+\pi)+\sin(x+2\pi).
\]
再令 \(t=x+\pi\)，得
\[
f(x+\pi)=f(x)+\sin(x+\pi).
\]
代入上式：
\[
f(x+2\pi)=f(x)+\sin(x+\pi)+\sin(x+2\pi).
\]
利用诱导公式
\[
\sin(x+\pi)=-\sin x,\qquad \sin(x+2\pi)=\sin x,
\]
可得
\[
\sin(x+\pi)+\sin(x+2\pi)=0.
\]
因此
\[
f(x+2\pi)=f(x).
\]
故 \(2\pi\) 是 \(f(x)\) 的一个周期。

\textbf{合法性检查：}

题设说明 \(f(x)\) 在全体实数上满足关系式，所以 \(x+\pi\)、\(x+2\pi\) 都在定义域内，递推代入合法。题目只要求证明 \(f(x)\) 以 \(2\pi\) 为周期，不要求证明 \(2\pi\) 是最小正周期。

\textbf{考场满分写法：}

由题设，
\[
f(x+2\pi)=f(x+\pi)+\sin(x+2\pi).
\]
又
\[
f(x+\pi)=f(x)+\sin(x+\pi).
\]
故
\[
f(x+2\pi)=f(x)+\sin(x+\pi)+\sin(x+2\pi)
=f(x)-\sin x+\sin x=f(x).
\]
所以 \(f(x)\) 以 \(2\pi\) 为周期。

\textbf{答案：}

\[
\boxed{f(x+2\pi)=f(x),\quad \text{故 } 2\pi\text{ 是 }f(x)\text{ 的一个周期}.}
\]
\end{solutionBox}

\begin{knowledgeBox}
\textbf{核心知识点：}

由递推关系证明周期性时，核心是“多次套用同一个等式，把目标 \(f(x+T)\) 拉回 \(f(x)\)”。本题的抵消来自
\[
\sin(x+\pi)+\sin(x+2\pi)=0.
\]

\textbf{可迁移方法：}

遇到
\[
f(x)=f(x-a)+\varphi(x)
\]
这类题，要尝试连续迭代：
\[
f(x+2a),\quad f(x+a),\quad f(x),
\]
观察新增项能否抵消。如果新增项是三角函数，优先使用诱导公式。
\end{knowledgeBox}

\begin{mistakeBox}
\textbf{易错点：}

\begin{itemize}
  \item \textbf{只套一次等式：} 只写 \(f(x+2\pi)=f(x+\pi)+\sin(x+2\pi)\) 还没有回到 \(f(x)\)。
  \item \textbf{诱导公式符号错：} \(\sin(x+\pi)=-\sin x\)，\(\sin(x+2\pi)=\sin x\)。
  \item \textbf{误以为要证明最小正周期：} 本题只需证明 \(2\pi\) 是一个周期。
\end{itemize}

\textbf{复盘提醒：}

周期证明题先写目标 \(f(x+T)=f(x)\)。题设若是递推式，就连续套用，直到所有新增项抵消。
\end{mistakeBox}

\studySubsection{calc-01-recover-function-from-composition}{由复合形式反求函数表达式}

\problemAnchor{MATH1-CALC-0001}
\begin{problemBox}
\begin{problemMeta}
\textbf{题目编号：} \problemRef{MATH1-CALC-0001} \quad
\textbf{索引：} \problemIndexRef{MATH1-CALC-0001} \quad
\textbf{来源：} 用户图片 / 教材例 1.1

\textbf{题型：} 已知 \(f(\varphi(u))\) 反求 \(f(x)\) \quad
\textbf{难度：} 基础

\textbf{知识点：} \knowledgeIndexRef[MATH1-KN-CALC-0002]{整体换元}，函数表示，\(u+\frac{1}{u}\) 的值域
\end{problemMeta}

\textbf{题目：}

设
\[
f\left(u+\frac{1}{u}\right)=\frac{u+u^3}{1+u^4},
\]
则当 \(x\ge 2\) 时，求 \(f(x)\)。

\textbf{题面校注：} 用户图片中分子首项疑似为 \(-u\)，但图片所给参考答案 \(\frac{x}{x^2-2}\) 对应的是分子为 \(u+u^3\) 的版本。若分子确认为 \(-u+u^3\)，则还需要指定 \(u\) 的取值分支，否则 \(f\left(u+\frac{1}{u}\right)\) 不能唯一确定。
\end{problemBox}

\begin{solutionBox}
\textbf{思路入口：}

目标是求 \(f(x)\)，已知的是 \(f\left(u+\frac{1}{u}\right)\)。所以不要把左右两边的 \(x\) 混用，应先把原式中的自变量记作 \(u\)，再把
\[
t=u+\frac{1}{u}
\]
看成一个整体。

\textbf{详细解答：}

由题设
\[
f\left(u+\frac{1}{u}\right)=\frac{u+u^3}{1+u^4}.
\]
为了把右边也化成 \(u+\frac{1}{u}\) 的形式，对分子、分母同时除以 \(u^2\)。这里要求 \(u\ne 0\)，而 \(u+\frac{1}{u}\) 本身也只在 \(u\ne 0\) 时有意义：
\[
\frac{u+u^3}{1+u^4}
=\frac{\frac{1}{u}+u}{\frac{1}{u^2}+u^2}.
\]
又因为
\[
u^2+\frac{1}{u^2}
=\left(u+\frac{1}{u}\right)^2-2,
\]
所以
\[
f\left(u+\frac{1}{u}\right)
=\frac{u+\frac{1}{u}}{\left(u+\frac{1}{u}\right)^2-2}.
\]
令
\[
t=u+\frac{1}{u},
\]
则
\[
f(t)=\frac{t}{t^2-2}.
\]
题目要求 \(x\ge 2\) 时的 \(f(x)\)，于是把 \(t\) 换回 \(x\)，得到
\[
f(x)=\frac{x}{x^2-2}\qquad (x\ge 2).
\]

\textbf{补充解法：反解参数后消元}

也可以先把最终自变量记为 \(t\)，令
\[
t=u+\frac{1}{u}\qquad (t\ge 2).
\]
这等价于
\[
u^2-tu+1=0.
\]
由此可得
\[
1+u^2=tu.
\]
于是
\[
u+u^3=u(1+u^2)=tu^2.
\]
另一方面，由
\[
u^2+\frac{1}{u^2}=\left(u+\frac{1}{u}\right)^2-2=t^2-2
\]
可得
\[
1+u^4=u^2(t^2-2).
\]
所以
\[
f(t)=\frac{u+u^3}{1+u^4}
=\frac{tu^2}{u^2(t^2-2)}
=\frac{t}{t^2-2}.
\]
因此
\[
f(x)=\frac{x}{x^2-2}\qquad (x\ge 2).
\]

\textbf{关于显式解出 \(u\) 再代回：}

也可以直接解二次方程
\[
u^2-tu+1=0,
\]
得到
\[
u=\frac{t\pm\sqrt{t^2-4}}{2}.
\]
但这里不能随意写成单值的 \(u=\theta(t)\)，因为当 \(t>2\) 时有两个分支，它们互为倒数。若将 \(u\) 换成 \(\frac{1}{u}\)，则
\[
\frac{\frac{1}{u}+\frac{1}{u^3}}{1+\frac{1}{u^4}}
=\frac{u^3+u}{u^4+1}
=\frac{u+u^3}{1+u^4},
\]
所以本题两个分支代回右端得到同一个值，显式解根代回是可行的。只是直接代入根式会使运算变繁，考场上不如先用方程关系消元。

\textbf{方法比较：} 本法本质上是先由 \(t=u+\frac{1}{u}\) 建立 \(u\) 与 \(t\) 的方程关系，再用方程消去 \(u\)。它能做，但比整体换元多了一层消元；考场上优先使用前一种整体改写法。若要显式写 \(u=\theta(t)\)，必须先说明分支对右端函数值没有影响。

\textbf{答案：}

\[
\boxed{f(x)=\frac{x}{x^2-2}\quad (x\ge 2)}.
\]
\end{solutionBox}

\begin{knowledgeBox}
\textbf{核心知识点：}

已知 \(f(\varphi(u))\) 时，关键不是先解出 \(u\)，而是把右端改写成只含 \(\varphi(u)\) 的形式。本题中
\[
u^2+\frac{1}{u^2}=\left(u+\frac{1}{u}\right)^2-2
\]
是主要变形入口。

\textbf{可迁移方法：}

遇到 \(u+\frac{1}{u}\)、\(u-\frac{1}{u}\)、\(u^2+\frac{1}{u^2}\) 同时出现的题，优先考虑整体换元。若题目限定 \(x\ge 2\)，通常是在提示 \(x\) 可以作为 \(u+\frac{1}{u}\) 的函数值。
\end{knowledgeBox}

\begin{mistakeBox}
\textbf{易错点：}

\begin{itemize}
  \item 把题设中的参数 \(u\) 和最终要求的自变量 \(x\) 混成同一个变量。
  \item 只想到解方程 \(x=u+\frac{1}{u}\)，却忽略右端可以直接化成整体变量。
  \item 忘记检查题面符号：若分子是 \(-u+u^3\)，则右端为 \(\frac{u-\frac{1}{u}}{u^2+\frac{1}{u^2}}\)，不能直接得到参考答案。
\end{itemize}

\textbf{复盘提醒：}

看到 \(f\left(u+\frac{1}{u}\right)\) 这类题，先重命名参数，再问自己：右端能不能通过除以某个 \(u^k\) 造出 \(u+\frac{1}{u}\)？
\end{mistakeBox}

\studySubsection{calc-01-reciprocal-equation-method}{互换 \(x\) 与 \(\frac{1}{x}\) 建立方程组求函数}

\problemAnchor{MATH1-CALC-0002}
\begin{problemBox}
\begin{problemMeta}
\textbf{题目编号：} \problemRef{MATH1-CALC-0002} \quad
\textbf{索引：} \problemIndexRef{MATH1-CALC-0002} \quad
\textbf{来源：} 用户图片 / 教材例 1.2

\textbf{题型：} 函数关系式中含 \(f(x)\) 与 \(f\left(\frac{1}{x}\right)\)，通过互换自变量建立方程组 \quad
\textbf{难度：} 基础

\textbf{知识点：} 定义域，\knowledgeIndexRef[MATH1-KN-CALC-0005]{倒数代换}，二元一次方程组，根式化简
\end{problemMeta}

\textbf{题目：}

设函数 \(f(x)\) 的定义域为 \((0,+\infty)\)，且满足
\[
2f(x)+x^2f\left(\frac{1}{x}\right)=\frac{x^2+2x}{\sqrt{1+x^2}},
\]
则求 \(f(x)\)。
\end{problemBox}

\begin{solutionBox}
\textbf{思路入口：}

题设中同时出现 \(f(x)\) 与 \(f\left(\frac{1}{x}\right)\)。要求的是 \(f(x)\)，但题设只给出一个关系式，未知量却可看作两个：
\[
f(x),\qquad f\left(\frac{1}{x}\right).
\]
因此自然的入口是把题设中的 \(x\) 换成 \(\frac{1}{x}\)，再得到第二个方程，联立消去 \(f\left(\frac{1}{x}\right)\)。

\textbf{详细解答：}

由题设，对任意 \(x>0\)，有
\[
2f(x)+x^2f\left(\frac{1}{x}\right)=\frac{x^2+2x}{\sqrt{1+x^2}}.
\qquad\text{(1)}
\]
因为定义域为 \((0,+\infty)\)，当 \(x>0\) 时也有 \(\frac{1}{x}>0\)，所以可以在题设中把 \(x\) 换为 \(\frac{1}{x}\)，得到
\[
2f\left(\frac{1}{x}\right)+\frac{1}{x^2}f(x)
=\frac{\frac{1}{x^2}+\frac{2}{x}}{\sqrt{1+\frac{1}{x^2}}}.
\]
右端需要化简。由于 \(x>0\)，
\[
\sqrt{1+\frac{1}{x^2}}
=\sqrt{\frac{x^2+1}{x^2}}
=\frac{\sqrt{1+x^2}}{x}.
\]
因此
\[
\frac{\frac{1}{x^2}+\frac{2}{x}}{\sqrt{1+\frac{1}{x^2}}}
=\frac{\frac{1+2x}{x^2}}{\frac{\sqrt{1+x^2}}{x}}
=\frac{1+2x}{x\sqrt{1+x^2}}.
\]
于是得到第二个方程
\[
2f\left(\frac{1}{x}\right)+\frac{1}{x^2}f(x)
=\frac{1+2x}{x\sqrt{1+x^2}}.
\qquad\text{(2)}
\]
为了消去 \(f\left(\frac{1}{x}\right)\)，将 \((2)\) 乘以 \(x^2\)，得
\[
2x^2f\left(\frac{1}{x}\right)+f(x)
=\frac{x(1+2x)}{\sqrt{1+x^2}}.
\qquad\text{(3)}
\]
再将 \((1)\) 乘以 \(2\)，得
\[
4f(x)+2x^2f\left(\frac{1}{x}\right)
=\frac{2x^2+4x}{\sqrt{1+x^2}}.
\qquad\text{(4)}
\]
用 \((4)-(3)\)，左端中 \(f\left(\frac{1}{x}\right)\) 被消去：
\[
3f(x)
=\frac{2x^2+4x-x-2x^2}{\sqrt{1+x^2}}
=\frac{3x}{\sqrt{1+x^2}}.
\]
所以
\[
f(x)=\frac{x}{\sqrt{1+x^2}}\qquad (x>0).
\]

\textbf{答案：}

\[
\boxed{f(x)=\frac{x}{\sqrt{1+x^2}}\quad (x>0)}.
\]
\end{solutionBox}

\begin{knowledgeBox}
\textbf{核心知识点：}

当函数关系式中同时出现 \(f(x)\) 和 \(f\left(\frac{1}{x}\right)\) 时，常用办法是把 \(x\) 换成 \(\frac{1}{x}\)，得到关于两个未知量的方程组：
\[
f(x),\qquad f\left(\frac{1}{x}\right).
\]
然后通过加减消元求出目标函数。

\textbf{可迁移方法：}

类似地，若关系式中出现 \(f(x)\) 与 \(f(-x)\)，常用 \(x\mapsto -x\)；若出现 \(f(x)\) 与 \(f(a-x)\)，常用 \(x\mapsto a-x\)。核心不是代换本身，而是构造“同一组未知函数值”的方程组。
\end{knowledgeBox}

\begin{mistakeBox}
\textbf{易错点：}

\begin{itemize}
  \item 忽略定义域：本题能把 \(x\) 换成 \(\frac{1}{x}\)，依赖于 \(x>0\) 时 \(\frac{1}{x}>0\)。
  \item 根式化简漏掉正负：\(\sqrt{x^2}=|x|\)，但本题 \(x>0\)，所以 \(\sqrt{x^2}=x\)。
  \item 消元倍数选错：应让 \(f\left(\frac{1}{x}\right)\) 前的系数相同，再相减。
\end{itemize}

\textbf{复盘提醒：}

看到 \(f(x)\) 与 \(f\left(\frac{1}{x}\right)\) 同时出现，先写两个未知量，再做 \(x\mapsto \frac{1}{x}\) 的对偶代换，最后联立消元。
\end{mistakeBox}

% BEGIN calculus-foundation:calc-01-pages
\studySubsection{calc-01-function-basics}{函数、定义域与基本初等函数}

\begin{knowledgeBox}
本页补齐函数基础域的稳定节点。已有的复合函数、反函数、隐函数、有界性、单调性、奇偶性和周期性详细页面继续作为唯一正文，不在此重复。
\end{knowledgeBox}

\knowledgeAnchor[MATH1-KN-CALC-0001]{函数与对应关系}
\begin{knowledgeBox}
\textbf{定义或结论：} 函数是定义域 $D$ 到数集的单值映射；相同函数必须定义域相同且对应法则相同。

\begin{itemize}
  \item \textbf{直觉与必会动作：} 会从“每个自变量唯一对应函数值”判断是否为函数；区分函数表达式与函数本身。
  \item \textbf{成立条件：} 定义域非空；“相同解析式”不必代表相同函数，须同时比较定义域。
  \item \textbf{典型用法与识别信号：} 判断两个函数是否相同；由关系式识别函数；识别信号为 “同一表达式”“是否为函数”“函数相同”。
  \item \textbf{关系与迁移：} 前置知识为 集合、初等代数；可迁移到 与极限存在、连续性、隐函数结合。
  \item \textbf{常见误区：} 只比较解析式；忽略定义域或多值关系。迁移到新题时先复核定义域、趋近方向和所用定理的条件。
\end{itemize}
\end{knowledgeBox}

\knowledgeAnchor[MATH1-KN-CALC-0007]{基本初等函数定义域}
\begin{knowledgeBox}
\textbf{定义或结论：} 定义域由所有组成运算共同约束：分母 $\ne0$，偶次根式被开方数 $\ge0$，$\ln u$ 要求 $u>0$，$\arcsin u,\arccos u$ 要求 $|u|\le1$。

\begin{itemize}
  \item \textbf{直觉与必会动作：} 熟练处理分母非零、偶次根号内非负、对数真数正、反三角函数自变量范围。
  \item \textbf{成立条件：} 多个条件取交集；含参数时需讨论可行集是否为空。
  \item \textbf{典型用法与识别信号：} 求函数定义域；含参数定义域；识别信号为 “定义域”“有意义”“所有实数”。
  \item \textbf{关系与迁移：} 前置知识为 不等式、集合交集；可迁移到 复合函数、极限、积分区间、幂级数收敛域。
  \item \textbf{常见误区：} 把根式条件写成严格大于；忘记对数真数严格为正；漏分母。迁移到新题时先复核定义域、趋近方向和所用定理的条件。
\end{itemize}
\end{knowledgeBox}

\knowledgeAnchor[MATH1-KN-CALC-0061]{复合函数定义域}
\begin{knowledgeBox}
\textbf{定义或结论：} $f(g(x))$ 的定义域为 $\{x:x\in D_g,\ g(x)\in D_f\}$。

\begin{itemize}
  \item \textbf{直觉与必会动作：} 先令内层函数有定义，再要求其值落入外层函数定义域。
  \item \textbf{成立条件：} 不能仅把外层自变量符号机械替换；若题目只给抽象定义域，必须做集合逆像。
  \item \textbf{典型用法与识别信号：} 已知 $D_f,D_g$ 求 $D_{f\circ g}$；识别信号为 “$f(x)$定义域为…，求$f(g(x))$”。
  \item \textbf{关系与迁移：} 前置知识为 基本定义域、集合；可迁移到 与反函数、函数方程、连续性结合。
  \item \textbf{常见误区：} 只求 $g(x)$ 的定义域；把 $D_f$ 当成 $x$ 的范围。迁移到新题时先复核定义域、趋近方向和所用定理的条件。
\end{itemize}
\end{knowledgeBox}

\knowledgeAnchor[MATH1-KN-CALC-0062]{复合函数性质}
\begin{knowledgeBox}
\textbf{定义或结论：} 复合单调性取决于内外层单调方向；若 $g$ 的值域落入 $f$ 的单调区间，才可组合判断。

\begin{itemize}
  \item \textbf{直觉与必会动作：} 会判断复合函数的奇偶、单调、周期与有界性。
  \item \textbf{成立条件：} 必须限定内层值域；外层不是全域单调时不能直接套“同增异减”。
  \item \textbf{典型用法与识别信号：} 抽象函数选择题；复合函数单调性；识别信号为 “$f$单调，$g$单调”“$f(g(x))$”。
  \item \textbf{关系与迁移：} 前置知识为 单调性、定义域；可迁移到 反函数、极限、导数、隐函数。
  \item \textbf{常见误区：} 忽略内层值域跨越外层不同单调区间。迁移到新题时先复核定义域、趋近方向和所用定理的条件。
\end{itemize}
\end{knowledgeBox}

\knowledgeAnchor[MATH1-KN-CALC-0014]{分段函数}
\begin{knowledgeBox}
\textbf{定义或结论：} 分段函数的性质必须在各段与连接点分别检查；连续需左右极限等于函数值，可导需左右导数相等。

\begin{itemize}
  \item \textbf{直觉与必会动作：} 会分别研究各段，并在分界点用极限、函数值、左右导数拼接。
  \item \textbf{成立条件：} 分界点可能属于某一段或单独定义；先看定义。
  \item \textbf{典型用法与识别信号：} 求参数使连续/可导；判断极值或拐点；识别信号为 “分段”“参数a”“在x=0连续/可导”。
  \item \textbf{关系与迁移：} 前置知识为 左右极限；可迁移到 极限、连续、导数定义、中值定理。
  \item \textbf{常见误区：} 只比较两段表达式值；忘记函数值或左右导数。迁移到新题时先复核定义域、趋近方向和所用定理的条件。
\end{itemize}
\end{knowledgeBox}

\knowledgeAnchor[MATH1-KN-CALC-0063]{指数、对数、幂函数}
\begin{knowledgeBox}
\textbf{定义或结论：} $a^x$（$a>0,a\ne1$）、$\log_a x$（$x>0$）；实幂 $x^\alpha$ 的定义域依 $\alpha$ 而异，考研常在正邻域使用 $x^\alpha=e^{\alpha\ln x}$。

\begin{itemize}
  \item \textbf{直觉与必会动作：} 掌握定义域、单调性、图像、基本恒等式及参数影响。
  \item \textbf{成立条件：} 做对数化时底数与真数须合法；含 $x^x$ 等表达式通常要求 $x>0$ 或局部单侧讨论。
  \item \textbf{典型用法与识别信号：} 图像性质；指数型极限；求导积分；识别信号为 “$f(x)^{g(x)}$”“取对数”。
  \item \textbf{关系与迁移：} 前置知识为 初等代数；可迁移到 泰勒、幂指函数、微分方程。
  \item \textbf{常见误区：} 忽略实幂定义域；对负底数直接取对数。迁移到新题时先复核定义域、趋近方向和所用定理的条件。
\end{itemize}
\end{knowledgeBox}

\knowledgeAnchor[MATH1-KN-CALC-0064]{三角与反三角函数}
\begin{knowledgeBox}
\textbf{定义或结论：} 熟练使用 $\sin^2x+\cos^2x=1$、和差化积、倍角/半角，以及反三角函数主值范围。

\begin{itemize}
  \item \textbf{直觉与必会动作：} 掌握定义域、值域、周期、奇偶、单调区间和基本恒等变换。
  \item \textbf{成立条件：} 反三角函数复合恒等式须检查主值区间；例如 $\arcsin(\sin x)=x$ 仅在主值区间成立。
  \item \textbf{典型用法与识别信号：} 基本计算、积分换元、极限；识别信号为 出现三角有理式、反三角复合。
  \item \textbf{关系与迁移：} 前置知识为 三角恒等式；可迁移到 傅里叶级数、极坐标、参数方程。
  \item \textbf{常见误区：} 把反函数恒等式无条件全域使用；周期判断错误。迁移到新题时先复核定义域、趋近方向和所用定理的条件。
\end{itemize}
\end{knowledgeBox}

\knowledgeAnchor[MATH1-KN-CALC-0022]{平移、伸缩、对称与复合图像}
\begin{knowledgeBox}
\textbf{定义或结论：} 横向变换作用在自变量，纵向变换作用在函数值；$f(|x|)$ 由右半图像偶延拓，$|f(x)|$ 将负值翻到上方。

\begin{itemize}
  \item \textbf{直觉与必会动作：} 能由基本图像快速得到 $f(ax+b)$、$af(x)+b$、$|f(x)|$、$f(|x|)$。
  \item \textbf{成立条件：} 变换顺序按表达式层级执行；横向伸缩系数与直觉相反。
  \item \textbf{典型用法与识别信号：} 识图、区域描述、绝对值函数性质；识别信号为 “图像”“围成区域”“绝对值”。
  \item \textbf{关系与迁移：} 前置知识为 基本初等函数图像；可迁移到 积分区域、单调极值、反函数图像。
  \item \textbf{常见误区：} 把横向平移符号或伸缩倍数写反。迁移到新题时先复核定义域、趋近方向和所用定理的条件。
\end{itemize}
\end{knowledgeBox}

\knowledgeAnchor[MATH1-KN-CALC-0310]{定义域与复合函数}
\begin{knowledgeBox}
\textbf{题型识别：} 表达式含分母、根式、对数、反三角或抽象复合；常以选择/填空出现；识别信号为 “有意义”“定义域为”“已知f的定义域，求f(g(x))”。

\begin{itemize}
  \item \textbf{标准流程：} ①列内层自身约束；②列内层值落入外层定义域；③取交集；④含参数时按临界点分段；⑤检查集合非空。
  \item \textbf{方法选择：} 抽象复合优先集合逆像；具体表达式优先逐层约束；可替代路线包括 画数轴、不等式符号表、集合逆像。
  \item \textbf{合法性检查：} 每一步先核对定义域、极限或定理条件；特别防止 只看最外层；把外层定义域直接当x范围；漏严格不等式。
  \item \textbf{考场步骤：} 先写识别出的结构与必要条件，再按标准流程逐步化简，最后回代、筛根或复核单侧情形。
  \item \textbf{迁移方向：} 含参数、嵌套根式、反函数定义域、幂级数复合收敛域；可与 函数性质、极限、幂级数 交错训练。
\end{itemize}
\end{knowledgeBox}

\studySubsection{calc-01-function-limits}{函数极限的定义、运算与存在性}

\begin{knowledgeBox}
极限计算按“先判趋近过程与型态，再选合法工具，最后复核条件”的顺序组织；七种未定式与泰勒展开沿用既有稳定节点。
\end{knowledgeBox}

\knowledgeAnchor[MATH1-KN-CALC-0073]{有限点函数极限的定义}
\begin{knowledgeBox}
\textbf{定义或结论：} $\lim_{x\to x_0}f(x)=A$ 指对任意 $\varepsilon>0$，存在 $\delta>0$，当 $0<|x-x_0|<\delta$ 时 $|f(x)-A|<\varepsilon$。

\begin{itemize}
  \item \textbf{直觉与必会动作：} 理解去心邻域与函数值无关；会写左右极限定义。
  \item \textbf{成立条件：} 不要求 $f(x_0)$ 有定义；只研究去心邻域。
  \item \textbf{典型用法与识别信号：} 定义判断；极限存在与点值关系辨析；识别信号为 “极限存在但函数未定义”“去心邻域”。
  \item \textbf{关系与迁移：} 前置知识为 数列极限量词；可迁移到 连续性、可导性、间断点。
  \item \textbf{常见误区：} 把 $x=x_0$ 包含进极限定义；把极限等同函数值。迁移到新题时先复核定义域、趋近方向和所用定理的条件。
\end{itemize}
\end{knowledgeBox}

\knowledgeAnchor[MATH1-KN-CALC-0074]{左右极限与双侧极限}
\begin{knowledgeBox}
\textbf{定义或结论：} $\lim_{x\to x_0}f(x)$ 存在且等于 $A$ 当且仅当 $f(x_0^-)=f(x_0^+)=A$。

\begin{itemize}
  \item \textbf{直觉与必会动作：} 会分别计算并用“左右相等”判断双侧极限。
  \item \textbf{成立条件：} 端点问题通常只讨论定义域内的单侧连续；趋于无穷无左右之分。
  \item \textbf{典型用法与识别信号：} 分段函数极限；含绝对值、符号函数；识别信号为 “$x\to0$”“分段”“绝对值”。
  \item \textbf{关系与迁移：} 前置知识为 函数极限；可迁移到 连续、可导、变上限积分。
  \item \textbf{常见误区：} 只算一侧；左右都发散但方向不同却称极限为无穷。迁移到新题时先复核定义域、趋近方向和所用定理的条件。
\end{itemize}
\end{knowledgeBox}

\knowledgeAnchor[MATH1-KN-CALC-0075]{无穷处与无穷函数极限}
\begin{knowledgeBox}
\textbf{定义或结论：} $x\to\infty$ 描述自变量离开任意有界区间；$f(x)\to\infty$ 表示函数值最终超过任意给定正数。

\begin{itemize}
  \item \textbf{直觉与必会动作：} 会处理 $x\to\infty$、$x\to-\infty$ 与 $f(x)\to\infty$ 的定义和运算。
  \item \textbf{成立条件：} “无穷大”描述趋势，不是实数；$+\infty$ 与 $-\infty$ 方向要区分。
  \item \textbf{典型用法与识别信号：} 有理、根式、指数对数在无穷处极限；识别信号为 “当x趋于无穷”“无界”。
  \item \textbf{关系与迁移：} 前置知识为 增长阶比较；可迁移到 渐近线、反常积分、级数。
  \item \textbf{常见误区：} 把 $\infty$ 当数参与减法；忽略正负无穷方向。迁移到新题时先复核定义域、趋近方向和所用定理的条件。
\end{itemize}
\end{knowledgeBox}

\knowledgeAnchor[MATH1-KN-CALC-0076]{函数极限四则运算与复合极限}
\begin{knowledgeBox}
\textbf{定义或结论：} 有限极限满足四则运算；若 $g(x)\to u_0$ 且 $f$ 在 $u_0$ 连续，则 $f(g(x))\to f(u_0)$。

\begin{itemize}
  \item \textbf{直觉与必会动作：} 会在条件满足时直接代入；识别不能拆分的未定式。
  \item \textbf{成立条件：} 商法则要求分母极限非零；复合换元须保证内层值落在外层定义域且趋近方式合法。
  \item \textbf{典型用法与识别信号：} 直接代入型极限；复合极限；识别信号为 “代入得非零有限值”。
  \item \textbf{关系与迁移：} 前置知识为 初等函数连续性；可迁移到 连续性、等价无穷小、泰勒。
  \item \textbf{常见误区：} 对 $0/0$ 强行拆；内层只从一侧趋近却套双侧外极限。迁移到新题时先复核定义域、趋近方向和所用定理的条件。
\end{itemize}
\end{knowledgeBox}

\knowledgeAnchor[MATH1-KN-CALC-0077]{函数极限的夹逼准则}
\begin{knowledgeBox}
\textbf{定义或结论：} 在某去心邻域有 $g(x)\le f(x)\le h(x)$，且两端同趋于 $A$，则 $f(x)\to A$。

\begin{itemize}
  \item \textbf{直觉与必会动作：} 会利用绝对值、积分界、三角界构造夹逼。
  \item \textbf{成立条件：} 邻域内不等式必须成立；夹逼对象可以是绝对值形式 $|f(x)-A|\le\varphi(x)\to0$。
  \item \textbf{典型用法与识别信号：} 振荡极限；积分平均值极限；识别信号为 “有界因子×无穷小”“积分区间缩短”。
  \item \textbf{关系与迁移：} 前置知识为 不等式；可迁移到 定积分、微分中值定理、级数余项。
  \item \textbf{常见误区：} 只写“显然夹逼”不写上下界；估计不趋零。迁移到新题时先复核定义域、趋近方向和所用定理的条件。
\end{itemize}
\end{knowledgeBox}

\knowledgeAnchor[MATH1-KN-CALC-0078]{函数极限的序列判据（海涅定理）}
\begin{knowledgeBox}
\textbf{定义或结论：} $\lim_{x\to x_0}f(x)=A$ 当且仅当对任意 $x_n\to x_0$ 且 $x_n\ne x_0$，有 $f(x_n)\to A$。

\begin{itemize}
  \item \textbf{直觉与必会动作：} 理解函数极限存在等价于任意趋近点列的函数值都趋同一极限。
  \item \textbf{成立条件：} “任意点列”是关键；证明极限必须控制所有点列，而否定一个候选极限只需构造一列反例。
  \item \textbf{见证数列法：} 若能找到两列都趋于 \(x_0\)，但函数值分别趋向不同极限，或一列有界而另一列无界，则函数在 \(x_0\) 处没有有限极限。它是把邻域中的振荡“锁定”成可计算数列的通用方法。
  \item \textbf{典型用法与识别信号：} 证明极限不存在；区分有界、无界与无穷大量；识别信号为 “快速振荡”“沿不同路径/序列”“某些点取零而另一些点很大”。
  \item \textbf{关系与迁移：} 前置知识为 数列极限；可迁移到 多元极限路径法、连续性逻辑。
  \item \textbf{常见误区：} 用一个点列收敛来证明函数极限；只检查有限条路径后宣称存在。迁移到新题时先复核定义域、趋近方向和所用定理的条件。
\end{itemize}
\end{knowledgeBox}

\knowledgeAnchor[MATH1-KN-CALC-0079]{无穷小、无穷大及互为倒数}
\begin{knowledgeBox}
\textbf{定义或结论：} 若 $\alpha(x)\to0$，称其为无穷小；若 $|f(x)|\to\infty$，称为无穷大。非零无穷小的倒数为无穷大，反之亦然。无穷大量必在每个去心邻域无界，但无界量未必是无穷大量。

\begin{itemize}
  \item \textbf{量词差别：} 在 \(x_0\) 的每个去心邻域无界是
  \[
    \forall\delta>0\ \forall M>0\ \exists x:
    0<|x-x_0|<\delta,\quad |f(x)|>M;
  \]
  而 \(|f(x)|\to\infty\) 是
  \[
    \forall M>0\ \exists\delta>0\ \forall x:
    0<|x-x_0|<\delta\Longrightarrow |f(x)|>M.
  \]
  前者只要求邻域中“总能找到大值点”，后者要求“所有充分靠近的点都取大值”。
  \item \textbf{双见证数列判据：} 构造 \(x_n\to x_0\) 且 \(|f(x_n)|\to\infty\)，即可证明无界；再构造 \(y_n\to x_0\) 且 \(|f(y_n)|\) 有界，即可否定无穷大量。两列合用便得到“无界但不是无穷大量”。
  \item \textbf{振幅乘振荡因子：} 对 \(F(x)=A(x)B(x)\)，若 \(|A(x)|\to\infty\) 而 \(B\) 振荡，先选“峰值列”使 \(B(x_n)\to c\ne0\)，再选“消去列”使 \(A(y_n)B(y_n)\) 保持有界。前者负责证明无界，后者负责否定无穷大量。
  \item \textbf{直觉与必会动作：} 会在指定过程中判断；区分“某些点任意大”和“所有邻近点最终都大”；区分无穷小量与常数0。
  \item \textbf{成立条件：} 必须明确趋近过程；倒数结论要求在去心邻域非零。
  \item \textbf{典型用法与识别信号：} 判断无穷小阶；区分无界量与无穷大量；处理发散振幅乘周期或振荡因子；识别信号为 “无穷小量”“无穷大量”“振荡且振幅增大”。
  \item \textbf{关系与迁移：} 前置知识为 极限定义；可迁移到 等价替换、级数必要条件、反常积分。
  \item \textbf{常见误区：} 由 \(|B(x)|\le C\) 和 \(|A(x)|\to\infty\) 就断言 \(|A(x)B(x)|\to\infty\)；一个趋于无穷的上界不能推出被估计量趋于无穷。迁移到新题时先分别寻找峰值列与消去列。
\end{itemize}
\end{knowledgeBox}

\knowledgeAnchor[MATH1-KN-CALC-0080]{同阶、高阶、低阶与等价无穷小}
\begin{knowledgeBox}
\textbf{定义或结论：} 若 $\alpha/\beta\to c\ne0$ 为同阶；趋0则 $\alpha=o(\beta)$；趋1则 $\alpha\sim\beta$。

\begin{itemize}
  \item \textbf{直觉与必会动作：} 会用比值极限分类，并掌握 $o,O,\sim$ 的含义。
  \item \textbf{成立条件：} $\beta$ 在去心邻域需非零；等价是同阶的特殊情形，不是相等。
  \item \textbf{典型用法与识别信号：} 判断无穷小阶；求参数使同阶/等价；识别信号为 “几阶无穷小”“同阶”“等价”。
  \item \textbf{关系与迁移：} 前置知识为 极限四则；可迁移到 泰勒展开、导数定义、级数比较。
  \item \textbf{常见误区：} 把同阶误作等价；忽略系数；把 $O$ 与 $o$ 混用。迁移到新题时先复核定义域、趋近方向和所用定理的条件。
\end{itemize}
\end{knowledgeBox}

\knowledgeAnchor[MATH1-KN-CALC-0081]{常用等价无穷小}
\begin{knowledgeBox}
\textbf{定义或结论：} $\sin x\sim x$，$\tan x\sim x$，$1-\cos x\sim x^2/2$，$e^x-1\sim x$，$\ln(1+x)\sim x$，$(1+x)^\alpha-1\sim\alpha x$，$\arcsin x\sim x$，$\arctan x\sim x$。

\begin{itemize}
  \item \textbf{直觉与必会动作：} 熟记并能推导 $x\to0$ 时三角、指数、对数、根式等价式。
  \item \textbf{成立条件：} 仅在相应自变量趋0时使用；复合后要求内层趋0；参数导致首项消失时需升阶。
  \item \textbf{典型用法与识别信号：} 乘商型极限快速计算；识别信号为 “趋于0的基本函数差”。
  \item \textbf{关系与迁移：} 前置知识为 基本极限；可迁移到 泰勒、积分等价、微分方程局部分析。
  \item \textbf{常见误区：} 漏系数1/2或α；内层不趋0；在加减式中直接逐项替换造成消项错误。迁移到新题时先复核定义域、趋近方向和所用定理的条件。
\end{itemize}
\end{knowledgeBox}

\knowledgeAnchor[MATH1-KN-CALC-0082]{等价无穷小替换规则}
\begin{knowledgeBox}
\textbf{定义或结论：} 若 $\alpha\sim\tilde\alpha$，在不发生危险消去的乘商结构中可替换；对 $\alpha-\beta$，首项可能相消，必须比较更高阶。

\begin{itemize}
  \item \textbf{直觉与必会动作：} 理解等价替换主要安全用于乘积和商；加减结构先提取或展开。
  \item \textbf{成立条件：} 替换后整体相对误差须可控；分母不能把误差放大。
  \item \textbf{典型用法与识别信号：} 复杂极限方法选择；判断错误解法；识别信号为 “两个近似量相减”“首项抵消”。
  \item \textbf{关系与迁移：} 前置知识为 等价无穷小；可迁移到 泰勒、洛必达、积分中值定理。
  \item \textbf{常见误区：} 把 $\sin x-x$ 直接替换为0；在和差中机械替换。迁移到新题时先复核定义域、趋近方向和所用定理的条件。
\end{itemize}
\end{knowledgeBox}

\knowledgeAnchor[MATH1-KN-CALC-0083]{第一重要极限}
\begin{knowledgeBox}
\textbf{定义或结论：} $\lim_{x\to0}\frac{\sin x}{x}=1$；由此得到 $\tan x/x\to1$、$\arcsin x/x\to1$ 等。

\begin{itemize}
  \item \textbf{直觉与必会动作：} 掌握 $\lim_{x\to0}\sin x/x=1$ 及其变形。
  \item \textbf{成立条件：} 角度使用弧度；复合内层须趋0。
  \item \textbf{典型用法与识别信号：} 三角极限、参数极限；识别信号为 “sin(小量)/小量”。
  \item \textbf{关系与迁移：} 前置知识为 三角函数；可迁移到 等价无穷小、导数公式。
  \item \textbf{常见误区：} 把角度制代入；分子分母的小量不对应。迁移到新题时先复核定义域、趋近方向和所用定理的条件。
\end{itemize}
\end{knowledgeBox}

\knowledgeAnchor[MATH1-KN-CALC-0084]{第二重要极限与 $e$}
\begin{knowledgeBox}
\textbf{定义或结论：} 若 $u\to0$ 且 $v\to\infty$，并有 $uv\to A$，在底数为正时 $(1+u)^v\to e^A$。

\begin{itemize}
  \item \textbf{直觉与必会动作：} 掌握 $(1+1/n)^n\to e$、$(1+x)^{1/x}\to e$ 及一般化。
  \item \textbf{成立条件：} 底数在邻域内必须为正以定义实幂；一般幂指型应取对数。
  \item \textbf{典型用法与识别信号：} $1^\infty$ 型极限；参数极限；识别信号为 “底数→1，指数→∞”。
  \item \textbf{关系与迁移：} 前置知识为 指数对数；可迁移到 对数化、泰勒、连续复利模型。
  \item \textbf{常见误区：} 只看成 $1^\infty=1$；忽略底数正性；漏掉高阶影响。迁移到新题时先复核定义域、趋近方向和所用定理的条件。
\end{itemize}
\end{knowledgeBox}

\knowledgeAnchor[MATH1-KN-CALC-0085]{代数化简、因式分解与有理化}
\begin{knowledgeBox}
\textbf{定义或结论：} 极限允许在去心邻域内用等价表达式替换；因式分解后可约去不恒为零的公因子。

\begin{itemize}
  \item \textbf{直觉与必会动作：} 会针对 $0/0$ 消去公因子、通分、共轭有理化。
  \item \textbf{成立条件：} 约分只要求去心邻域非零，不要求点上有定义；有理化需同时乘共轭。
  \item \textbf{典型用法与识别信号：} 根式极限、有理式极限；识别信号为 “根号差”“高次多项式差”。
  \item \textbf{关系与迁移：} 前置知识为 因式分解；可迁移到 泰勒、洛必达前的预处理。
  \item \textbf{常见误区：} 约去可能在邻域内多次为零的因子而不说明；代数错误。迁移到新题时先复核定义域、趋近方向和所用定理的条件。
\end{itemize}
\end{knowledgeBox}

\knowledgeAnchor[MATH1-KN-CALC-0086]{幂指型极限的对数化}
\begin{knowledgeBox}
\textbf{定义或结论：} 令 $y=f(x)^{g(x)}>0$，则 $\ln y=g(x)\ln f(x)$；先求右侧极限 $L$，再得 $y\to e^L$。

\begin{itemize}
  \item \textbf{直觉与必会动作：} 统一处理 $0^0,1^\infty,\infty^0$ 型。
  \item \textbf{成立条件：} 底数需在所考邻域为正；若对数极限为 $\pm\infty$，相应判断0或无穷。
  \item \textbf{典型用法与识别信号：} $x^x$、复合幂函数极限；识别信号为 “变量底数的变量指数”。
  \item \textbf{关系与迁移：} 前置知识为 对数函数；可迁移到 泰勒、洛必达、指数连续性。
  \item \textbf{常见误区：} 对非正底数直接取对数；求完 $\ln y$ 忘记指数还原。迁移到新题时先复核定义域、趋近方向和所用定理的条件。
\end{itemize}
\end{knowledgeBox}

\knowledgeAnchor[MATH1-KN-CALC-0312]{常规极限方法选择}
\begin{knowledgeBox}
\textbf{题型识别：} 极限可呈0/0、∞/∞、0·∞、∞−∞或直接代入型；识别信号为 代入后出现未定式；含基本小量、根式、积分或幂指。

\begin{itemize}
  \item \textbf{标准流程：} ①先判趋近过程与型态；②代数化简；③乘商看等价无穷小；④高阶消去用Taylor；⑤商型可比较洛必达；⑥振荡/积分用夹逼；⑦最后核验条件。
  \item \textbf{方法选择：} 低阶结构用等价；需要系数/高阶用Taylor；连续可导商且求导变简单用洛必达；可替代路线包括 变量代换、导数定义、积分中值定理。
  \item \textbf{合法性检查：} 每一步先核对定义域、极限或定理条件；特别防止 见0/0即洛必达；加减中机械等价；展开阶数不足。
  \item \textbf{考场步骤：} 先写识别出的结构与必要条件，再按标准流程逐步化简，最后回代、筛根或复核单侧情形。
  \item \textbf{迁移方向：} 含参数使极限有限、含绝对值单侧、积分平均值、数列和式；可与 Taylor、导数、积分、级数 交错训练。
\end{itemize}
\end{knowledgeBox}

\knowledgeAnchor[MATH1-KN-CALC-0313]{幂指型极限}
\begin{knowledgeBox}
\textbf{题型识别：} 底数和指数均变化，常见 $1^\infty,0^0,\infty^0$；识别信号为 “$f(x)^{g(x)}$”“$x^x$”“底数趋1”。

\begin{itemize}
  \item \textbf{标准流程：} ①确认底数在邻域为正；②令y为原式并取ln；③求 $g\ln f$ 的极限；④用等价/Taylor/洛必达；⑤指数还原；⑥处理±∞。
  \item \textbf{方法选择：} 结构可化 $(1+u)^v$ 且uv易求时用重要极限；否则统一对数化；可替代路线包括 改写为第二重要极限；直接Taylor指数。
  \item \textbf{合法性检查：} 每一步先核对定义域、极限或定理条件；特别防止 把 $1^\infty$ 当1；非正底数强取ln；忘还原。
  \item \textbf{考场步骤：} 先写识别出的结构与必要条件，再按标准流程逐步化简，最后回代、筛根或复核单侧情形。
  \item \textbf{迁移方向：} 参数底数、指数积分、两层指数；可与 极限、对数、Taylor 交错训练。
\end{itemize}
\end{knowledgeBox}

\studySubsection{calc-01-continuity-theorems}{连续性、间断点与闭区间定理}

\begin{knowledgeBox}
连续性问题先核对点值与左右极限，再使用闭区间定理。边界知识只保留检索入口，不进入本地复习安排。
\end{knowledgeBox}

\knowledgeAnchor[MATH1-KN-CALC-0087]{函数在一点连续}
\begin{knowledgeBox}
\textbf{定义或结论：} $f$ 在 $x_0$ 连续当且仅当 $\lim_{x\to x_0}f(x)=f(x_0)$，即 $\Delta x\to0$ 时 $\Delta y\to0$。

\begin{itemize}
  \item \textbf{直觉与必会动作：} 会用极限等于函数值判断；理解连续包含有定义、极限存在、二者相等。
  \item \textbf{成立条件：} 双侧连续要求两侧均在定义域内；端点通常用单侧连续。
  \item \textbf{典型用法与识别信号：} 判断连续；求参数补定义；识别信号为 “连续”“补充定义”“参数”。
  \item \textbf{关系与迁移：} 前置知识为 函数极限；可迁移到 可导性、积分上限函数、闭区间定理。
  \item \textbf{常见误区：} 漏检查函数值；极限存在就称连续。迁移到新题时先复核定义域、趋近方向和所用定理的条件。
\end{itemize}
\end{knowledgeBox}

\knowledgeAnchor[MATH1-KN-CALC-0088]{左连续、右连续与区间连续}
\begin{knowledgeBox}
\textbf{定义或结论：} 左连续：$f(x_0^-)=f(x_0)$；右连续类似。闭区间连续指内部连续且端点相应单侧连续。

\begin{itemize}
  \item \textbf{直觉与必会动作：} 会在端点和分段点使用单侧连续。
  \item \textbf{成立条件：} 端点外侧若不在定义域，无需讨论；分段函数边界归属要看定义。
  \item \textbf{典型用法与识别信号：} 分段参数题；闭区间定理条件核查；识别信号为 “在[a,b]连续”“右连续”。
  \item \textbf{关系与迁移：} 前置知识为 左右极限；可迁移到 最值、介值、积分存在性。
  \item \textbf{常见误区：} 要求端点双侧连续；忽略分段点属于哪段。迁移到新题时先复核定义域、趋近方向和所用定理的条件。
\end{itemize}
\end{knowledgeBox}

\knowledgeAnchor[MATH1-KN-CALC-0089]{间断点分类}
\begin{knowledgeBox}
\textbf{定义或结论：} 第一类间断点左右极限均存在：相等为可去，不等为跳跃；至少一侧极限不存在或无穷为第二类。

\begin{itemize}
  \item \textbf{直觉与必会动作：} 会区分可去、跳跃、无穷和振荡间断。
  \item \textbf{成立条件：} “函数值错误或未定义”可形成可去间断；分类依据左右极限，不依据图像是否断开。
  \item \textbf{典型用法与识别信号：} 判间断点及类型；求参数消除间断；识别信号为 “间断点类型”“可去”。
  \item \textbf{关系与迁移：} 前置知识为 左右极限；可迁移到 反常积分、分段函数、周期延拓。
  \item \textbf{常见误区：} 把无穷间断归为跳跃；只看点值。迁移到新题时先复核定义域、趋近方向和所用定理的条件。
\end{itemize}
\end{knowledgeBox}

\knowledgeAnchor[MATH1-KN-CALC-0090]{初等函数与复合函数连续性}
\begin{knowledgeBox}
\textbf{定义或结论：} 基本初等函数在其定义域内连续；有限次四则运算与合法复合保持连续。

\begin{itemize}
  \item \textbf{直觉与必会动作：} 会在定义域内直接用连续性代入求极限。
  \item \textbf{成立条件：} 商要求分母非零；复合要求内层值落在外层连续点。
  \item \textbf{典型用法与识别信号：} 直接代入型极限；证明函数连续；识别信号为 “初等函数”“定义域内”。
  \item \textbf{关系与迁移：} 前置知识为 基本函数；可迁移到 变上限积分、多元连续性。
  \item \textbf{常见误区：} 忽略分母为零或外层边界点。迁移到新题时先复核定义域、趋近方向和所用定理的条件。
\end{itemize}
\end{knowledgeBox}

\knowledgeAnchor[MATH1-KN-CALC-0091]{有界性与最大最小值定理}
\begin{knowledgeBox}
\textbf{定义或结论：} 若 $f\in C[a,b]$，则存在 $x_m,x_M\in[a,b]$ 使 $f(x_m)=m,f(x_M)=M$。

\begin{itemize}
  \item \textbf{直觉与必会动作：} 理解闭区间连续函数必有界且能取到最大、最小值。
  \item \textbf{成立条件：} 闭区间与连续缺一不可；开区间或仅有界不保证取到最值。
  \item \textbf{典型用法与识别信号：} 判断定理适用；证明存在最值；识别信号为 “闭区间连续”“一定存在”。
  \item \textbf{关系与迁移：} 前置知识为 连续性；可迁移到 全局最值、积分估计、多元紧集性质。
  \item \textbf{常见误区：} 把局部极值当全局最值；漏端点比较。迁移到新题时先复核定义域、趋近方向和所用定理的条件。
\end{itemize}
\end{knowledgeBox}

\knowledgeAnchor[MATH1-KN-CALC-0092]{零点定理与介值定理}
\begin{knowledgeBox}
\textbf{定义或结论：} 若 $f\in C[a,b]$ 且 $f(a)f(b)<0$，则存在 $\xi\in(a,b)$ 使 $f(\xi)=0$；介值定理更一般。

\begin{itemize}
  \item \textbf{直觉与必会动作：} 会通过端点异号证明至少一个零点；会证明函数取某中间值。
  \item \textbf{成立条件：} 只保证存在，不保证唯一；端点同号也不能推出无根。
  \item \textbf{典型用法与识别信号：} 证明方程有根；根区间定位；识别信号为 “证明至少有一根”“端点异号”。
  \item \textbf{关系与迁移：} 前置知识为 连续性、不等式；可迁移到 单调性证明唯一根、中值定理、参数方程。
  \item \textbf{常见误区：} 把存在性与唯一性混淆；函数不连续仍套零点定理。迁移到新题时先复核定义域、趋近方向和所用定理的条件。
\end{itemize}
\end{knowledgeBox}

\knowledgeAnchor[MATH1-KN-CALC-0093]{一致连续性（教材扩展）}
\begin{knowledgeBox}
\textbf{边界定位：} 一致连续要求同一个 $\delta$ 能同时服务于定义域内所有点；闭区间上的连续函数具有这种更强的整体控制。

\textbf{学习边界：} 本仓库把它作为边界知识：用于理解闭区间连续性的延伸，不展开证明，也不进入本地复习安排。

\textbf{误区与迁移：} 常见误区是把逐点连续直接当成一致连续，或把闭区间结论无条件搬到无界区间。
\end{knowledgeBox}

\knowledgeAnchor[MATH1-KN-CALC-0094]{柯西收敛准则与 Stolz 定理}
\begin{knowledgeBox}
\textbf{边界定位：} 柯西收敛准则与 Stolz--Ces\`aro 定理可作为数列极限的拓展视角。

\textbf{学习边界：} 本仓库只保留名称、用途以及它们与夹逼、单调有界方法的关系；规范解答优先使用本讲已整理的基础工具，不展开拓展证明，也不进入本地复习安排。

\textbf{误区与迁移：} 常见误区是只会套拓展定理，却没有先证明适用条件或没有掌握基本判据。
\end{knowledgeBox}

\knowledgeAnchor[MATH1-KN-CALC-0311]{分段点连续与可导参数}
\begin{knowledgeBox}
\textbf{题型识别：} 分界点两侧表达式不同，要求连续、可导或二阶可导并求参数；识别信号为 “在x=0处连续/可导”“确定a,b”。

\begin{itemize}
  \item \textbf{标准流程：} ①算左右极限和点值，先解连续条件；②算左右导数，解可导条件；③二阶问题再逐侧求导；④回代核验。
  \item \textbf{方法选择：} 表达式复杂或含绝对值时定义法更稳；解析展开明显时用系数匹配；可替代路线包括 直接用导数定义；Taylor匹配低阶系数。
  \item \textbf{合法性检查：} 每一步先核对定义域、极限或定理条件；特别防止 跳过连续直接比导数；把某段公式在分界点直接代入；漏左右。
  \item \textbf{考场步骤：} 先写识别出的结构与必要条件，再按标准流程逐步化简，最后回代、筛根或复核单侧情形。
  \item \textbf{迁移方向：} 积分定义函数、参数方程拼接、连续但不可导反例；可与 极限、导数定义、Taylor 交错训练。
\end{itemize}
\end{knowledgeBox}
% END calculus-foundation:calc-01-pages
